\setlength\beforechapskip{-1em}
\chapter{Laplacians}
In  \Cref{chap:gaussianFields} the fractional Gaussian fields will be defined using the eigenvalues and eigenfunctions of the Dirichlet Laplacian. We will find that the Dirichlet Laplacian is an extension of the Kigami Laplacian. But to construct the Kigami Laplacian we need the discrete Laplacian as we will prove that the Kigami Laplacian is the pointwise limit of the discrete Laplacian. We will therefore start by introducing the discrete Laplacian, then the Kigami Laplacian, culminating with the Dirichlet Laplacian. Therefore consider the discrete Laplacian on the Vicsek set.

\medskip
\begin{defi}\label{discreteLaplacian}
	Let $\ell(V_m)$ be the set of functions $\func{f}{V_m}{\ree}$. Then for any $f \in \ell(V_m)$  the \emph{discrete Laplacian} on $V_m$ is defined as
	\begin{equation}
		\Delta_m f(p) = 3^{m} \frac{1}{N_p}\sum_{q \in V_{m,p}} (f(q) - f(p)), \quad p \in V_m \setminus V_0
	\end{equation}
	where $$N_p = \begin{cases}
		5^{-m}, & \text{if } \ \# V_{m,p} = 4 \\
		2\cd 5^{-(m+1)}, & \text{if } \ \# V_{m,p} = 2\\
		5^{-(m+1)},  & \text{if } \ \# V_{m,p} = 1.
	\end{cases}$$
\end{defi}

\medskip
One is able to describe the discrete Laplacian of a function $f$ as a matrix called the Laplacian matrix.
\medskip
\begin{defi}\label{discreteMatrix}
	 Let $p_1, \ldots p_{N_m}$ be the vertices in $V_m$ then the \emph{Laplacian matrix} $L_m$ is given by 
	\begin{equation}
		(L_m)_{i,j} = \begin{cases}
			\#V_{m,p_i} \cd 3^{m} \cd\frac{1}{N_{p_i}}, & \text{if } i= j \\
			-1 \cd 3^{m} \cd \frac{1}{N_{p_i}}, & \text{if } i \neq j \text{ and } p_j \in V_{m,p_i} \\
			0, & \text{otherwise}
		\end{cases}
	\end{equation}
	which satisfies:
	\begin{equation}
		L_mF = \begin{pmatrix}
			\Delta_m f(x_1)\\
			\vdots \\
			\Delta_{m}f(x_{N_m})
		\end{pmatrix}, \quad \text{where } F = \begin{pmatrix}
			f(x_1)\\
			\vdots \\
			f(x_{N_m})
		\end{pmatrix}.
	\end{equation}
\end{defi}

\section{Graph Energies}
In this section we will be following the arguments given by Strichartz in \cite[Chapter 1]{strichartz}. Many of the proofs are going to rely on the graph energy. Thus, we will first define and prove some properties of the graph energies.

\medskip
\begin{defi}
	Given a finite connected graph $G_m$ and it vertices $V_m$ then for functions $u, v\colon V_m \to \ree$ we define the \emph{graph energy} to be
	\begin{equation}\label{graphEnergy}
		E_m(u,v) = \sum_{x \sim y} \left(u(x) - u(y) \right)(v(x)-v(y)).
	\end{equation}
\end{defi}
\textbf{Notation:} We will write $E_m(u) = E_m(u,u)$.

\medskip
 Since we want to work on the entire Vicsek set $K$, we will want to create an extension of a function $u\colon V_m \to \ree$, such that it is defined on $V_{m+n}$. This is done in the following way:

\medskip
\begin{defi}
	Let $u\colon V_m \to \ree$. Then an \emph{extension of} $u$ is a function $u'\colon V_{m+n} \to \ree$ where $u' \mid_{V_m} = u$.
\end{defi}
  One quickly notices that there must be a function that minimizes the graph energy of a given graph. We denote such a function by $\tilde{u}$ and call it a \emph{harmonic extension}.

\medskip
\begin{defi}
	Let $\tilde{u}\colon V_{m+n} \to \ree$ with $\tilde{u} \mid_{V_m} = u$ being a \emph{harmonic extension} if \\ $E_{m+n}(\tilde{u}) \leq E_{m+n}(u')$ for all $u'\colon V_{m+n} \to \ree$ with $u' \mid _{V_m} = u$.
\end{defi}

 In some cases we have that $E_{m+n}(\tilde{u}) = r^{n} E_{m}(u)$ for some $r \in (0,1)$. But instead of normalizing the values on $E_m$ we will normalize $E_{m+n}$, thus  getting
\begin{equation}
	r^{-n} E_{m+n}(\tilde{u}) = E_m(u), \quad \text{and} \quad r^{-n} E_{m+n}(u') \geq E_m(u)
\end{equation}
for every extension $u'$ of $u$. The constant $r$ is called the \emph{renormalization constant}. From the equality above we can define the following.

\medskip
\begin{defi}\label{renormalizedGraphEnergies}
	The \emph{renormalized graph energies} are defined as
	\begin{equation}
		\eee_{m}(u) = r^{-m}E_m(u).
	\end{equation}
	Note that the sequence $\cbr{\eee_{m}(u)}_{m \geq 0}$ is a non-decreasing sequence. 
\end{defi}

\medskip
Based on this definition we see that for a harmonic extension $\tilde{u}$ of $u$ we have that 
\begin{equation}
	\mathcal{E}_{m+1}(\tilde{u}) = r^{-(m+1)}E_{m+1}(\tilde{u}) = r^{-m}E_{m}(u) = \mathcal{E}_{m}(u).
\end{equation}
Extending this to the bilinear form we have, for harmonic extension $\tilde{u}, \tilde{v}$ of $u, v$ that
\begin{equation}
	\mathcal{E}_{m+1}(\tilde{u}, \tilde{v}) = \mathcal{E}_{m}(u,v).
\end{equation}

It also true for the bilinear form, that if $\tilde{u}$ is the harmonic extension of $u$ and $v'$ is any extension of $v$ then
\begin{equation}\label{bilinearFormRenormalized}
	\mathcal{E}_{m+1}(\tilde{u}, v') = \mathcal{E}_{m}(u,v).
\end{equation}
For further details see \cite{strichartz}.

Now that we have \Cref{renormalizedGraphEnergies}, it is possible to define a harmonic function.

\medskip
\begin{defi}
	A function $u$ is a \emph{harmonic function} if $u$ minimizes $\eee_{m}(u)$ for all $m$, given boundary values on $V_0$.
\end{defi}

The harmonic functions will be used extensively throughout this section. We start by finding the renormalization constant for the Vicsek set.

\subsection{The Renormalization Constant on the Vicsek set}

In the previous section we talked about a harmonic function, but without ever constructing one. We will therefore in this section show how to construct a harmonic function for the Vicsek set for any initial boundary values $u(V_0)$.

\medskip
\noi Given $V_0 = \cbr{q_1, q_2, q_3, q_4, q_5}$ and $u\colon V_0 \to \ree$ we have that
\begin{equation}
	E_0(u) = \left(u(q_1) - u(q_2)\right)^2 + \left(u(q_1) - u(q_3)\right)^2 + \left(u(q_1) - u(q_4)\right)^2 + \left(u(q_1) - u(q_5)\right)^2.
\end{equation} 
We are going to exploit the symmetry of the Vicsek set. Since all four "extremities" of the Vicsek set are identical, see \Cref{renormVicsekDiagram}, we will only look at one of the extremities, since the other extremities will be the same computations. So if we want to minimize $E_1(u')$, for $u'\colon V_1 \to \ree$, we see for the extremity between $q_1$ and $q_2$ that
\begin{multline}\label{E1Energy}
	\hspace{2em} E_1(u') = \left(u'(q_2) - u'(p_1)\right)^2 + \left(u'(p_1) - u'(p_3)\right)^2 + \left(u'(p_1) - u'(p_4)\right)^2 
	\\ + \left(u'(p_1) - u'(p_5)\right)^2 + \left(u'(p_4) - u'(q_1)\right)^2. \hspace{2em}
\end{multline}


\begin{figure}[!ht] \label{renormVicsekDiagram}
	\centering
	\tikzset{every picture/.style={line width=0.75pt}} %set default line width to 0.75pt        
	
	\begin{tikzpicture}[x=0.75pt,y=0.75pt,yscale=-1,xscale=1]
		%uncomment if require: \path (0,300); %set diagram left start at 0, and has height of 300
		
		%Straight Lines [id:da6304295946642043] 
		\draw    (200,120) -- (200,240) ;
		\draw [shift={(200,240)}, rotate = 90] [color={rgb, 255:red, 0; green, 0; blue, 0 }  ][fill={rgb, 255:red, 0; green, 0; blue, 0 }  ][line width=0.75]      (0, 0) circle [x radius= 1.34, y radius= 1.34]   ;
		\draw [shift={(200,120)}, rotate = 90] [color={rgb, 255:red, 0; green, 0; blue, 0 }  ][fill={rgb, 255:red, 0; green, 0; blue, 0 }  ][line width=0.75]      (0, 0) circle [x radius= 1.34, y radius= 1.34]   ;
		%Straight Lines [id:da15009868113581337] 
		\draw [color={rgb, 255:red, 0; green, 0; blue, 0 }  ,draw opacity=1 ]   (140,180) -- (260,180) ;
		\draw [shift={(260,180)}, rotate = 0] [color={rgb, 255:red, 0; green, 0; blue, 0 }  ,draw opacity=1 ][fill={rgb, 255:red, 0; green, 0; blue, 0 }  ,fill opacity=1 ][line width=0.75]      (0, 0) circle [x radius= 1.34, y radius= 1.34]   ;
		\draw [shift={(140,180)}, rotate = 0] [color={rgb, 255:red, 0; green, 0; blue, 0 }  ,draw opacity=1 ][fill={rgb, 255:red, 0; green, 0; blue, 0 }  ,fill opacity=1 ][line width=0.75]      (0, 0) circle [x radius= 1.34, y radius= 1.34]   ;
		%Straight Lines [id:da6360387946571586] 
		\draw    (160,160) -- (160,200) ;
		\draw [shift={(160,200)}, rotate = 90] [color={rgb, 255:red, 0; green, 0; blue, 0 }  ][fill={rgb, 255:red, 0; green, 0; blue, 0 }  ][line width=0.75]      (0, 0) circle [x radius= 1.34, y radius= 1.34]   ;
		\draw [shift={(160,160)}, rotate = 90] [color={rgb, 255:red, 0; green, 0; blue, 0 }  ][fill={rgb, 255:red, 0; green, 0; blue, 0 }  ][line width=0.75]      (0, 0) circle [x radius= 1.34, y radius= 1.34]   ;
		%Straight Lines [id:da967738390323848] 
		\draw    (240,160) -- (240,200) ;
		\draw [shift={(240,200)}, rotate = 90] [color={rgb, 255:red, 0; green, 0; blue, 0 }  ][fill={rgb, 255:red, 0; green, 0; blue, 0 }  ][line width=0.75]      (0, 0) circle [x radius= 1.34, y radius= 1.34]   ;
		\draw [shift={(240,160)}, rotate = 90] [color={rgb, 255:red, 0; green, 0; blue, 0 }  ][fill={rgb, 255:red, 0; green, 0; blue, 0 }  ][line width=0.75]      (0, 0) circle [x radius= 1.34, y radius= 1.34]   ;
		%Straight Lines [id:da22729733453474055] 
		\draw    (220,140) -- (180,140) ;
		\draw [shift={(180,140)}, rotate = 180] [color={rgb, 255:red, 0; green, 0; blue, 0 }  ][fill={rgb, 255:red, 0; green, 0; blue, 0 }  ][line width=0.75]      (0, 0) circle [x radius= 1.34, y radius= 1.34]   ;
		\draw [shift={(220,140)}, rotate = 180] [color={rgb, 255:red, 0; green, 0; blue, 0 }  ][fill={rgb, 255:red, 0; green, 0; blue, 0 }  ][line width=0.75]      (0, 0) circle [x radius= 1.34, y radius= 1.34]   ;
		%Straight Lines [id:da18095273720450344] 
		\draw    (170,180) -- (185.5,180) -- (190,180) ;
		%Straight Lines [id:da24264648223015173] 
		\draw    (220,220) -- (180,220) ;
		\draw [shift={(180,220)}, rotate = 180] [color={rgb, 255:red, 0; green, 0; blue, 0 }  ][fill={rgb, 255:red, 0; green, 0; blue, 0 }  ][line width=0.75]      (0, 0) circle [x radius= 1.34, y radius= 1.34]   ;
		\draw [shift={(220,220)}, rotate = 180] [color={rgb, 255:red, 0; green, 0; blue, 0 }  ][fill={rgb, 255:red, 0; green, 0; blue, 0 }  ][line width=0.75]      (0, 0) circle [x radius= 1.34, y radius= 1.34]   ;
		%Straight Lines [id:da6311736325856517] 
		\draw    (180,180) -- (160,180) ;
		\draw [shift={(160,180)}, rotate = 180] [color={rgb, 255:red, 0; green, 0; blue, 0 }  ][fill={rgb, 255:red, 0; green, 0; blue, 0 }  ][line width=0.75]      (0, 0) circle [x radius= 1.34, y radius= 1.34]   ;
		\draw [shift={(180,180)}, rotate = 180] [color={rgb, 255:red, 0; green, 0; blue, 0 }  ][fill={rgb, 255:red, 0; green, 0; blue, 0 }  ][line width=0.75]      (0, 0) circle [x radius= 1.34, y radius= 1.34]   ;
		%Straight Lines [id:da868237351757014] 
		\draw    (240,180) -- (220,180) ;
		\draw [shift={(220,180)}, rotate = 180] [color={rgb, 255:red, 0; green, 0; blue, 0 }  ][fill={rgb, 255:red, 0; green, 0; blue, 0 }  ][line width=0.75]      (0, 0) circle [x radius= 1.34, y radius= 1.34]   ;
		\draw [shift={(240,180)}, rotate = 180] [color={rgb, 255:red, 0; green, 0; blue, 0 }  ][fill={rgb, 255:red, 0; green, 0; blue, 0 }  ][line width=0.75]      (0, 0) circle [x radius= 1.34, y radius= 1.34]   ;
		%Straight Lines [id:da7980791265955824] 
		\draw    (200,220) -- (200,200) ;
		\draw [shift={(200,200)}, rotate = 270] [color={rgb, 255:red, 0; green, 0; blue, 0 }  ][fill={rgb, 255:red, 0; green, 0; blue, 0 }  ][line width=0.75]      (0, 0) circle [x radius= 1.34, y radius= 1.34]   ;
		\draw [shift={(200,220)}, rotate = 270] [color={rgb, 255:red, 0; green, 0; blue, 0 }  ][fill={rgb, 255:red, 0; green, 0; blue, 0 }  ][line width=0.75]      (0, 0) circle [x radius= 1.34, y radius= 1.34]   ;
		%Straight Lines [id:da7384137154814034] 
		\draw    (200,160) -- (200,140) ;
		\draw [shift={(200,140)}, rotate = 270] [color={rgb, 255:red, 0; green, 0; blue, 0 }  ][fill={rgb, 255:red, 0; green, 0; blue, 0 }  ][line width=0.75]      (0, 0) circle [x radius= 1.34, y radius= 1.34]   ;
		\draw [shift={(200,160)}, rotate = 270] [color={rgb, 255:red, 0; green, 0; blue, 0 }  ][fill={rgb, 255:red, 0; green, 0; blue, 0 }  ][line width=0.75]      (0, 0) circle [x radius= 1.34, y radius= 1.34]   ;
		%Straight Lines [id:da3701955076323855] 
		\draw    (200,200) -- (200,180) ;
		\draw [shift={(200,180)}, rotate = 270] [color={rgb, 255:red, 0; green, 0; blue, 0 }  ][fill={rgb, 255:red, 0; green, 0; blue, 0 }  ][line width=0.75]      (0, 0) circle [x radius= 1.34, y radius= 1.34]   ;
		
		% Text Node
		\draw (200,170.4) node [anchor=north west][inner sep=0.75pt]  [font=\scriptsize]  {$q_{1}$};
		% Text Node
		\draw (262,171.4) node [anchor=north west][inner sep=0.75pt]  [font=\scriptsize]  {$q_{2}$};
		% Text Node
		\draw (194,108) node [anchor=north west][inner sep=0.75pt]  [font=\scriptsize]  {$q_{5}$};
		% Text Node
		\draw (194,243.4) node [anchor=north west][inner sep=0.75pt]  [font=\scriptsize]  {$q_{3}$};
		% Text Node
		\draw (126,171.4) node [anchor=north west][inner sep=0.75pt]  [font=\scriptsize]  {$q_{4}$};
		% Text Node
		\draw (241,170.4) node [anchor=north west][inner sep=0.75pt]  [font=\scriptsize]  {$p_{1}$};
		% Text Node
		\draw (220,168.4) node [anchor=north west][inner sep=0.75pt]  [font=\scriptsize]  {$p_{4}$};
		% Text Node
		\draw (234,204) node [anchor=north west][inner sep=0.75pt]  [font=\scriptsize]  {$p_{3}$};
		% Text Node
		\draw (235,148) node [anchor=north west][inner sep=0.75pt]  [font=\scriptsize]  {$p_{5}$};
				
	\end{tikzpicture}
	\caption{Points on the Vicsek set $V_1$.}
\end{figure}
Next we are going to minimize the graph energies with regards to $u'(p_i)$ for $i = 0, 1, 2, 4$. Taking the derivative of the right side of \eqref{E1Energy} and solving for $0$ we get the following
\begin{align}
	4 u'(p_1) &= u'(q_2) + u'(p_3) + u'(p_4) + u'(p_5) \\
	2 u'(p_4) &= u'(q_1) + u'(p_1) \\
	u'(p_3) &= u'(p_1)\\
	u'(p_5) &= u'(p_1).
\end{align}
Solving these equations we get
\begin{align}
	u'(p_1) &=  \frac{2}{3} u'(q_2) + \frac{1}{3} u'(q_1)\\
	u'(p_3) &= \frac{2}{3} u'(q_2) + \frac{1}{3} u'(q_1) \\
	u'(p_5) &= \frac{2}{3} u'(q_2) + \frac{1}{3} u'(q_1) \\
	u'(p_4) &= \frac{2}{3} u'(q_1) + \frac{1}{3} u'(q_2).
\end{align}

Ergo what we have found is, that to minimize the value in the points of the Vicsek set, the value in the next point is going to be $2/3$ the value from the closest point and $1/3$ the value from the second closest point. 

If we calculate the graph energy we find that
\begin{equation}
	E_1(u') =\frac{1}{3}\sum_{i = 1}^{4}(u'(q_i) - u'(q_1))^2 = \frac{1}{3}E_0(u').
\end{equation}

Since the Vicsek set is a self-similar set, we have that going from $V_1$ to $V_2$ is the same as going from $V_m$ to $V_{m+1}$. Thus we can simply extend the above calculations repeating it $m$ times, in which case
\begin{equation}
	E_{m+1}(u') = \frac{1}{3}E_m(u').
\end{equation}

We then get our renormalization constant for the Vicsek set to be $r = \frac{1}{3}$. Altogether we find from \Cref{renormalizedGraphEnergies} that
\begin{equation}
	\mathcal{E}_{m}(u) = 3^{m} E_{m}(u).
\end{equation}


\section{The Kigami Laplacian}
With the renormalization constant found, we will return to finding the Kigami Laplacian. Recall from \Cref{renormalizedGraphEnergies} that we have defined a sequence $\mathcal{E}_m$ for functions on $V^{*}$. It will therefore make sense to define the following limit.

\medskip
\begin{defi}\label{renormalizedGraphEnergiesLimit} Let $u\colon V^{*} \to \ree$ then define the limit
	\begin{equation}
		\mathcal{E}(u, u) = \lim_{m \to \infty} \mathcal{E}_{m}(u, u),
	\end{equation}
	where we set
	\begin{equation}
		\dom \mathcal{E} = \cbr{u \in \ell(V_{*}) : \mathcal{E}(u) < \infty},\quad \text{and} \quad \dom_0 \mathcal{E} = \cbr{u \in \dom \mathcal{E} : u = 0,\text{ on } V_0 }.
	\end{equation}
\end{defi}

This definition only involves functions defined on $V^{*}$, but we find that functions in $\dom \mathcal{E}$ are uniformly continuous and therefore have a unique extension to $K$. We even have that $\dom \mathcal{E}$ is dense in $C(K)$.

\medskip
\begin{prop} Let $\dom \mathcal{E}$ be defined as in \Cref{renormalizedGraphEnergiesLimit}. Then the following holds
	\begin{enumerate}[(i)] 
		\item 	Let $u \in \dom \mathcal{E}$, then $u$ is uniformly continuous on $V^{*}$. \label{uniCont1}
		\item $\dom \mathcal{E}$ is dense in $C(K)$. \label{uniCont2}
	\end{enumerate}
\end{prop}
\begin{proof}
	We start with (\ref{uniCont1}). Let $u \in \dom \mathcal{E}$ then $\mathcal{E}_m(u) \leq \mathcal{E}(u)$ since $\mathcal{E}_m$ is non-decreasing in $m$. So for any $x,y \in V_m$ if $y \in V_{m,x}$ we have that
	\begin{equation}\label{uniFormContiousEstimate}
		r^{-m}(u(x) -u(y))^{2} \leq \mathcal{E}_m(u) \leq \mathcal{E}(u)
	\end{equation}
	since $r^{-m}(u(x) -u(y))^{2}$ is a term in $\mathcal{E}_m(u)$. Thus
	\begin{equation}
		\abs{u(x)-u(y)} \leq r^{\sfrac{m}{2}} \mathcal{E}(u)^{\sfrac{1}{2}}.
	\end{equation}
	
	Now consider a chain of points $x_m, x_{m+1}, \ldots, x_{m+k}$ where $x_{m+j} \in V_{m+j}$ and \\ $x_{m+j} \in V_{m+j+1, x_{m+j+1}}$. Then we have
	\begin{equation}
		\abs{u(x_m)-u(x_{m+k})} \leq r^{\sfrac{m}{2}}\left(1 + r^{\sfrac{1}{2}} + \ldots + r^{\sfrac{k}{2}}\right)\mathcal{E}(u)^{\sfrac{1}{2}} \leq \frac{r^{m}}{1-r^{\sfrac{1}{2}}} \mathcal{E}^{\sfrac{1}{2}}
	\end{equation}
	by using \eqref{uniFormContiousEstimate} and the triangle inequality multiple times. From the geometry of $K$ we see that $x, y \in V^{*}$ belong to the same or adjacent $\left(\fracConst_{i_1} \circ \ldots \circ \fracConst_{i_m}\right)(K)$. We can then connect $x$ to $y$ by one of such chains, so
	\begin{equation}
		\abs{u(x) - u(y)} \leq \frac{2r^{\sfrac{m}{2}}}{1-r^{\sfrac{1}{2}}} \mathcal{E}(u)^{\sfrac{1}{2}}
	\end{equation}
	giving a bound for all $x,y \in V^{*}$.
	
	For proof of (\ref{uniCont2}) we refer to \cite{strichartz}.
\end{proof}


We are now ready to define the Kigami Laplacian. We will denote this with $\Delta_{\mu}$ as it depends on the choice of measure $\mu$. In our case we choose $\mu$ to be given as in \Cref{normalziedHausdorffMeasure}.

\medskip
\begin{defi}
	Let $u \in \dom \mathcal{E}$ and let $f \in C(K)$, then $u \in \dom \Delta_{\mu}$ and $\Delta_{\mu} u =f$ if
	\begin{equation}
		\mathcal{E}(u,v) = \int_{K}fv \dd \mu, \quad \forall v \in \dom_{0}\mathcal{E}.
	\end{equation}
	We call $\Delta_{\mu}$ the \emph{Kigami Laplacian}.
\end{defi}

As mentioned before, we want $\Delta_{\mu}$ to be defined as the limit of $\Delta_m$. The constants $3^m$ and $5^m$ are not random, but are  chosen such that the Kigami Laplacian makes sense. But to prove the pointwise limit and actually calculate the constants, we will need a few lemmas and definitions. 

We will introduce a new notation for the discrete graph Laplacian without its constants. The reason for this will become apparent, since we in future proofs will see that the constants cancels each other out. We start by defining this graph Laplacian followed by a harmonic function which will be used extensively in the following proofs.

\medskip
\begin{defi}
	Let $\Delta_{m}$ be defined as in \Cref{discreteLaplacian}. We then define $\widetilde{\Delta}_m$ as
	\begin{equation}
		\widetilde{\Delta}_{m}f(p) = 3^{-m} N_p \Delta_mf(p) = \sum_{q \in V_{m, p}} f(q)-f(p).
	\end{equation}
\end{defi}

\begin{defi}
	Given $m \geq 0$ and a point $x \in V_m \backslash V_0$ then define the \emph{harmonic function} \\ $\psi_{x}^{(m)}\colon V^{*} \to [0,1]$ that for a $y \in V_m$ is $\psi_x^{(m)}(y) = \delta_{xy}$.
\end{defi}
Note that $\psi_x^{(m)} \in \dom \mathcal{E}$. Furthermore, it might seem odd that we only defined $\psi_x^{(m)}$ to have values on $V_m$. However the function is defined on the whole of $V^{*}$, where the values in $K \setminus V_m$ come from the fact that it is a harmonic extension. We even have that function $\psi_x^{(m)}$ can be used to describe $\widetilde{\Delta}_{m}$.
\medskip
\begin{prop}\label{pointwiseEqualityofpsi}
	Let $u \in \dom \mathcal{E}$ and $\Delta_{\mu}u = f$ then
	\begin{equation}
		r^{-m}\widetilde{\Delta}_{m}u(x) = \int_{K}f \psi_{x}^{(m)} \dd \mu.
	\end{equation}
\end{prop}
\begin{proof}
	We see that
	\begin{multline}
		\mathcal{E}_{m}(u, \psi_{x}^{(m)}) = r^{-m} \sum_{p \underset{m}{\sim} q} \left(u(q)-u(p)\right)\left(\psi_{x}^{(m)}(p) -\psi_{x}^{(m)}(q)\right) \\
		= r^{-m}\sum_{q \in V_{m,x}} \left(u(x)-u(q)\right) = r^{-m}\widetilde{\Delta}_{m}u(x),
	\end{multline}
	where the second equality stems from the fact that only edges containing $x$ are going to see change in $\psi_{x}^{(m)}$. By \eqref{bilinearFormRenormalized} we have that
	\begin{equation}
		r^{-m}\widetilde{\Delta}_{m}u(x) = \mathcal{E}_{m}(u, \psi_{x}^{(m)}) = \mathcal{E}(u, \psi_{x}^{(m)}) = \int_{K}f \psi_{x}^{(m)} \dd \mu
	\end{equation}
	as desired.
\end{proof}

\begin{lemma}\label{itIsZero}
	Let $v\colon V^{*} \to \ree$ be a harmonic function on the Vicsek set. If $v(x) = v(y_i) = 0$ for all $y_i \in V_{m,x}$ then $v(\omega) = 0$ for all $\omega \in \overline{T_x}$ where $T_x = \left( \bigcup_{k > 0} S^{k}(x) \right) \cup \left(\bigcup_{i} \{y_i\} \right)$.
\end{lemma}
\begin{proof}
	It is clear that $v(\omega) = 0$ for any $\omega \in T_x$, since the value of the points $S^{k}(x)$ with $k > 0$ depends on the value of $x$ and its neighbours. 
	
	So we need to check the points in $\overline{T_x} \setminus T_x$. Now since $T_x$ is dense in $\overline{T_x}$, then for any $\omega \in \overline{T_x}$ we have that there exists a sequence $x_n \in T_x$ such that $x_n \to \omega$. By the fact that $v$ is continuous on $K$ it is continuous on $\overline{T_x}$, in which case $\abs{v(x_n) - v(\omega)} \to 0$ as $n \to \infty$, but $v(x_n) = 0$ for all $x_n$, so $v(\omega) = 0$ as desired.  
\end{proof}

What \Cref{itIsZero} states is, if any two points in $V_m$ are neighbours, say $x$ and $y$, and $v(x)=v(y)=0$, then the function value at any point between $x$ and $y$ is also zero.

\medskip
\begin{lemma}\label{FmSmaller}
	Set $F_m = \supp \psi_{x}^{(m)}$, then
	\begin{equation}
		F_0 \supset F_1 \supset F_2 \supset F_3 \supset \dots.
	\end{equation}
	Further $F_m$ is contained in the open ball $B(x, \epsilon_m)$ where $\epsilon_m \to 0$ as $m \to \infty$.
\end{lemma}

\begin{proof}
	If we set $\Omega_m = \left\{z \in K : \psi_{x}^{(m)}(z) = \psi_{x}^{(m)}(y) = 0, \forall y \in V_{m,z}  \right\}$ then $F_m = K \setminus \left(\bigcup_{\alpha \in \Omega_m} \overline{B_\alpha} \right)$. It is clear that $\Omega_m \subset \Omega_{m+1}$ and that there are new points in $\Omega_{m+1}$ that are not just a point between some already existing points $x_1, x_2 \in \Omega_m$. We therefore have, given the construction of $F_m$, that $F_m \supset F_{m+1}$.
	
	The containment of $F_m$ in $B(x, \epsilon_m)$ follows from \Cref{itIsZero} and that $F_m \supset F_{m+1}$.
\end{proof}

We are now ready to prove the pointwise formula.

\medskip
\begin{theorem} \label{pointwiseFormula}
	Let $u \in \dom \Delta_{\mu}$. Then the pointwise formula
	\begin{equation}\label{pointwiseEquation}
		\Delta_{\mu} u(x) = \lim_{m \to \infty} r^{-m} \left(\int_{K} \psi_x^{(m)} \dd \mu \right)^{-1} \widetilde{\Delta}_m u(x)
	\end{equation}
	holds with the limit uniform across $V_* \setminus V_0$. Conversely, suppose $u$ is a continuous function and the right side of \eqref{pointwiseEquation} converges uniformly to a continuous function on $V_* \setminus V_0$. Then $u \in \dom \Delta_{\mu}$ and \eqref{pointwiseEquation} holds.
\end{theorem}

\begin{proof}
	Suppose $u \in \dom \Delta_{\mu}$ and $\Delta_{\mu} u=f$, then by \Cref{pointwiseEqualityofpsi} we have that 
	\begin{equation}
		r^{-m} \left(\int_{K} \psi^{(m)}_{x} \dd\mu \right)^{-1} \widetilde{\Delta}_m u(x) = \frac{\int_{K} f \psi_x^{(m)} \dd\mu}{\int_{K} \psi_x^{(m)} \dd\mu}
	\end{equation}
	for any $x \in V_m \setminus V_0$. The right side needs to converge uniformly, so let $\epsilon > 0$. Since $f(x)$ is continuous we can by \Cref{FmSmaller} choose an $m$ large enough such that for all $x, y \in F_m$ we have that $\abs{f(x) - f(y)} < \epsilon$, in which case
	\begin{align}
		\abs{\frac{\int_{K} f \psi_x^{(m)} \dd\mu}{\int_{K} \psi_x^{(m)} \dd\mu} - f(x)} &= \abs{\frac{\int_{K} f \psi_x^{(m)} \dd\mu}{\int_{K} \psi_x^{(m)} \dd\mu} - \frac{\int_{K} f(x) \psi_x^{(m)} \dd\mu}{\int_{K} \psi_x^{(m)} \dd\mu}} \\
		&= \abs{\frac{\int_{K} \left(f(y) -f(x) \right) \psi_x^{(m)} \dd\mu}{\int_{K} \psi_x^{(m)} \dd\mu}} \\
		&\leq \frac{\int_{F_m} \abs{f(y) -f(x)} \dd\mu}{\int_{K} \psi_x^{(m)} \dd\mu} \\
		&\leq \frac{\epsilon \int_{F_m} \dd\mu}{\int_{K} \psi_x^{(m)} \dd\mu} \\
		&= \epsilon,
	\end{align}
	therefore
	\begin{equation}
		\lim_{m \to \infty} r^{-m} \left(\int_{K} \psi^{(m)}_{x} \dd\mu \right)^{-1} \widetilde{\Delta}_m u(x) = \lim_{m \to \infty} \frac{\int_{K} f \psi_x^{(m)} \dd \mu}{\int_{K} \psi_x^{(m)} \dd \mu} = f(x) = \Delta_{\mu} u(x)
	\end{equation}
	proving the first statement.
	
	Next suppose that the expression $r^{m} \left(\int_{K} \psi_x^{(m)} \dd \mu \right)^{-1} \widetilde{\Delta}_m u(x)$ converges uniformly to a continuous function $f$. Then for any $v \in \dom_0 \mathcal{E}$ we have 
	\begin{equation}
		\mathcal{E}_m(u,v) = r^{-m} \sum_{x \underset{m}{\sim} y} \left( u(y) - u(x) \right)\left(v(y) - v(x) \right).
	\end{equation}
	
	If we collect all the terms that contain the factor $v(x)$ for a fixed $x \in V_m \setminus V_0$ then each $v(x)$ will be of the form $-v(x) \left(u(y) - u(x) \right)$ for each $y \underset{m}{\sim} x$. We therefore get
	\begin{align}
		\mathcal{E}_m(u,v) &= - r^{-m}  \sum_{x \in V_m \setminus V_0} v(x) \left( \sum_{x \underset{m}{\sim} y} \left( u(y) - u(x) \right) \right) \\
		&= -\sum_{x \in V_m \setminus V_0} v(x) r^{-m} \widetilde{\Delta}_{m}u(x).
	\end{align}
	While it seems we lose many terms by removing $v(y)$, they are instead contained in 
	$$\sum_{x \underset{m}{\sim} y} \left( u(y) - u(x) \right).$$ 
	Note that the sum is "reset" for every $x$ so we count both the relation $x \sim y$ and $y \sim x$.
	
	Now if we define $f_m(x) = r^{-m} \left(\int_{K} \psi_{x}^{(m)}(y) \dd\mu \right)^{-1} \widetilde{\Delta}_{m}u(x)$ then $f_m \to f$ uniformly, as $m \to \infty$, by assumption and we can rewrite the above as   
	\begin{align}
		\mathcal{E}_m(u,v) &= -\sum_{x \in V_m \setminus V_0} v(x) \left(\int_{K} \psi_{x}^{(m)}(y) \dd\mu \right) r^{-m} \left(\int_{K} \psi_{x}^{(m)}(y) \dd\mu \right)^{-1} \widetilde{\Delta}_{m}u(x) \\
		&= - \int_{K} \left( \sum_{x \in V_m \setminus V_0} v(x) f_m(x) \psi_x^{(m)}(y)\right) \dd \mu. \\
	\end{align}
	Since the convergence, in the next calculation, follows from a quite lengthy argument, we will show the calculation and then return to the convergence. We have that 
	\begin{multline}\label{theVeryLongOne}
		\mathcal{E}(u, v) = \lim_{m \to \infty} \mathcal{E}_m(u,v) \\
		= \lim_{m \to \infty}-\int_{K} \left( \sum_{x \in V_m \setminus V_0} v(x) f_m(x) \psi_x^{(m)}(y)\right) \dd \mu = - \int_{K} v(y)f(y) \dd \mu \quad
	\end{multline}
	therefore $u \in \text{dom}\Delta_{\mu}$. Now looking at the convergence in \eqref{theVeryLongOne} we find
	\begin{multline}
		\abs{\sum_{x \in V_m \setminus V_0} v(x) f_m(x) \psi_x^{(m)}(y) -v(y)f(y)} \\
		\leq \abs{\sum_{x \in V_m \setminus V_0} \left(v(x) f_m(x) - v(y)f(y) \right) \psi_x^{(m)}(y)} + \abs{\sum_{x \in V_m \setminus V_0} v(y) f(y) \psi_x^{(m)}(y) -v(y)f(y)}.
	\end{multline}
	Looking at each term individually we first see that
	\begin{multline}
		\abs{\sum_{x \in V_m \setminus V_0} \left(v(x) f_m(x) - v(y)f(y) \right) \psi_x^{(m)}(y)} \\ 
		\leq \sum_{x \in V_m \setminus V_0} \left(\abs{v(x) f_m(x) - v(y)f_m(y)} + \abs{v(y)f_m(y) - v(y)f(y)}\right) \psi_x^{(m)}(y).
	\end{multline}
	We will again look at each term starting with $\abs{v(y) f_m(y) - v(y)f(y)}$. Here
	$$\abs{v(y) f_m(y) - v(y)f(y)} = \abs{v(y)}\abs{f_m(y) - f(y)}$$
	and since $v$ is continuous on $K$, which is compact, there is a constant $C_1$ so $\abs{v(y)} \leq C_1$ for all $x \in K$. Because $f_m$ converges uniformly to $f$, we can choose $m_1$ such that $\abs{f_{m_1}(y)- f(y)} \leq \frac{\epsilon}{C_1}$ and the convergence of the term follows.
	
	Next is $\abs{v(x) f_m(x) - v(y)f_m(y)}$. Here we find
	\begin{equation}
		\abs{v(x) f_m(x) - v(y)f_m(y)} \leq \abs{v(x)}\abs{f_m(x)-f(y)} + \abs{f_m(y)}\abs{v(x)-v(y)}.
	\end{equation}
	First we look at $\abs{f_m(y)}\abs{v(x)-v(y)}$. Here $\abs{f_m(y)} \leq C_2$ and by \Cref{FmSmaller} and continuity of $v$ we can choose $m_2$ so $\abs{v(x)-v(y)} \leq \frac{\epsilon}{C_2}$ for all $x \in F_{m_2}$. The reason we can choose $m_2$ is because $\psi_{x}^{(m)}$ is multiplied onto all the terms restricting $v$ to $F_{m_2}$.
	
	\medskip
	Second we examine $\abs{v(x)}\abs{f_m(x)-f(y)}$. Again $\abs{v(x)} \leq C_1$, but
	\begin{equation}
		\abs{f_m(x)-f(y)} \leq \abs{f_m(x)-f(x)} + \abs{f(x) - f_m(y)} \leq \abs{f_m(x)-f(x)} + \abs{f(x) -f(y)} + \abs{f(y) - f_m(y)}
	\end{equation}
	First choose $m_3$ so $\abs{f_{m_3}(x)-f(x)} \leq \frac{\epsilon}{3 C_1}$. Next by continuity of $f$ choose $m_4$ so \\ $\abs{f(x) -f(y)} \leq \frac{\epsilon}{3 C_1}$ for all $x \in F_{m_4}$. Lastly choose $m_5$ so $\abs{f(y)-f_{m_5}(y)} \leq \frac{\epsilon}{3 C_1}$.
	
	So we now have for $m \geq \max \{m_1, m_2, \ldots, m_5\}$ that
	\begin{equation}
		\abs{\sum_{x \in V_m \setminus V_0} \left(v(x) f_m(x) - v(y)f(y) \right) \psi_x^{(m)}(y)} \leq \epsilon \sum_{x \in V_m \setminus V_0} \psi_x^{(m)}(y) \leq \epsilon.
	\end{equation}
	We are going to look at 
	\begin{equation}
		\abs{\sum_{x \in V_m \setminus V_0} v(y) f(y) \psi_x^{(m)}(y) -v(y)f(y)} = \abs{v(y)f(y)}\abs{\sum_{x \in V_m \setminus V_0} \psi_x^{(m)}(y) - 1}.
	\end{equation}
	and show that 
	\begin{equation}\label{psiConvergesToOne}
	\abs{\sum_{x \in V_m \setminus V_0} \psi_x^{(m)}(y) - 1} = 0, \quad \text{as } m \to \infty.
	\end{equation}
	If  $y \in V^*$, then choose $m_6$ so $y \in V_{m_6}$ and \eqref{psiConvergesToOne} holds.
	
	Next assume $y \in K \setminus V^{*}$, then since $V^{*}$ is dense in $K$ we can choose a sequence $(x_n)_{n \geq 1}$ where $x_n \to y$ as $n \to \infty$, and since $\psi_{x}^{(m)}$ is continuous on $K$ then $\psi_{x}^{(m)}(x_n) \to \psi_{x}^{(m)}(y)$ as $n \to \infty$. This gives us the following
	\begin{align}
		\abs{\sum_{x \in V_m \setminus V_0} \psi_x^{(m)}(y) - 1} \leq \abs{\sum_{x \in V_m \setminus V_0} \psi_x^{(m)}(y) - \sum_{x \in V_m \setminus V_0} \psi_x^{(m)}(x_n)} + \abs{\sum_{x \in V_m \setminus V_0} \psi_x^{(m)}(x_n) - 1}.
	\end{align}
	
	Now choose an $n_1$ so $\abs{\psi_{x}^{(m)}(x_{n_1}) - \psi_{x}^{(m)}(y)} \leq \epsilon$ and an $m_7$ so $x_{n_1} \in V_{m_7}$. Both can then be bounded by $m \geq \max \{n_1, m_7\}$. Thus we have the convergence of \eqref{theVeryLongOne}.
\end{proof}

Lastly we calculate the constants $\left(\int_{K} \psi_{x}^{(m)}(y) \dd \mu(y)\right)^{-1}$. 
\begin{lemma}
	\begin{equation}
		\int_{K} \psi_{x}^{m}(y) \dd \ \mu(y) = \begin{cases}
			5^{-m} & \text{if } \ \# V_{m,p} = 4 \\
			2\cd 5^{-(m+1)} & \text{if } \ \# V_{m,p} = 2 \\
			5^{-(m+1)}  & \text{if } \ \# V_{m,p} = 1 \: .
		\end{cases}
	\end{equation}
\end{lemma}

\begin{proof}
Recall from \cref{pointAsIntersectSet} that we can write a point $x$'s "address" with the infinite sequence $K_{i_1, i_2, \ldots} = (\fracConst_{i_1} \circ \fracConst_{i_2} \circ \ldots)(K)$. These sequences are not necessarily unique and we will have three cases. The first two cases are where $x$ has a unique sequence and $\# V_{m,p} \in \cbr{1,4}$. The last case is where $x$ has two sequences and $\# V_{m, p} = 2$. If we start with the two sequences, then we have that $x \in K_{i_1, \ldots, i_m} \cap K_{i_1, \ldots, i_{m'}}$ in which case from \Cref{itIsZero} it follows that $\supp \psi_{x}^{(m)} = K_{i_1, \ldots, i_{m+1}} \cup K_{i_1, \ldots, i_{m'+1}}$, for an appropriate choice of $i_{m+1}$ and $i_{m'+1}$. Now if $x\in K_{i_1, \ldots, i_m}$ and $\# V_{m,p} = 4$ we have that $\supp\psi_{x}^{(m)} = K_{i_1, \ldots, i_m}$, but if $\# V_{m,p} = 1$ then $\supp\psi_{x}^{(m)} = K_{i_1, \ldots, i_{m+1}}$ again for an appropriate $i_{m+1}$.

If we now look at an $\alpha \in \supp \psi_x^{(m)}$ it is going to have a sequence of the form $K_{i_1, \ldots i_m, \alpha_{j_1}, \alpha_{j_2} \ldots}$ or $K_{i_1, \ldots i_k, \alpha_{j_1}, \alpha_{j_2} \ldots}$ where $k \in \cbr{m+1, m'+1}$. We therefore have
\begin{equation}
	\int_{K} \psi_{x}^{(m)}(y) \dd \mu(y) = \int_{\supp \psi_{x}^{(m)}} \psi_{x}^{(m)}(y) \dd \mu(y) = \sum_{\alpha \in \supp \psi_{x}^{(m)}} \int_{\cbr{\alpha}} \psi_{x}^{(m)}(y) \dd \mu(y).
\end{equation}
 If $x$ has a unique sequence and $\# V_{m,p} = 4$ then
\begin{multline}
	\sum_{\alpha \in \supp \psi_{x}^{(m)}} \int_{\cbr{\alpha}} \psi_{x}^{(m)}(y) \dd \mu(y) = \sum_{\alpha \in \supp \psi_{x}^{(m)}} \mu\left((\fracConst_{i_1} \circ \ldots \circ \fracConst_{i_m} \circ \fracConst_{\alpha_{j_1}} \circ \fracConst_{\alpha_{j_2}} \ldots)(K) \right) \\
	= \mu\left(\bigcup_{(j_1, j_2, \ldots) \in \mathcal{I}} (\fracConst_{i_1} \circ \ldots \circ \fracConst_{i_m} \circ \fracConst_{\alpha_{j_1}} \circ \fracConst_{\alpha_{j_2}} \circ \ldots)(K)\right) = \mu \left((\fracConst_{i_1} \circ \ldots \circ \fracConst_{i_m})(K)\right) = 5^{-m}.
\end{multline}
By the almost exact same calculation we find that if $x$ has a unique sequence and $\# V_{m,p} = 1$ then 
\begin{equation}
	\sum_{\alpha \in \supp \psi_{x}^{(m)}} \int_{\cbr{\alpha}} \psi_{x}^{(m)}(y) \dd \mu(y) = 5^{-(m+1)}.
\end{equation}
Now if $x$ does not have a unique sequence then
\begin{align}
		\sum_{\alpha \in \supp \psi_{x}^{(m)}} \int_{\cbr{\alpha}} \psi_{x}^{(m)}(y) \dd \mu(y) &= \sum_{\alpha \in \supp \psi_{x}^{(m)}} \bigg[ \mu\left((\fracConst_{i_1} \circ \ldots \circ \fracConst_{i_m+1} \circ \fracConst_{\alpha_{j_1}} \circ \fracConst_{\alpha_{j_2}} \ldots)(K) \right) \\
		&\phantom{a} \qquad \qquad \qquad + \mu\left((\fracConst_{i_1} \circ \ldots \circ \fracConst_{i_m'+1} \circ \fracConst_{\alpha_{j_1}} \circ \fracConst_{\alpha_{j_2}} \ldots)(K) \right) \bigg] \\
		&= \mu\left(\bigcup_{(j_1, j_2, \ldots) \in \mathcal{I}} (\fracConst_{i_1} \circ \ldots \circ \fracConst_{i_m+1} \circ \fracConst_{\alpha_{j_1}} \circ \fracConst_{\alpha_{j_2}} \circ \ldots)(K)\right)\\
		&\phantom{a} \qquad \qquad \qquad + \mu\left(\bigcup_{(j_1, j_2, \ldots) \in \mathcal{I}} (\fracConst_{i_1} \circ \ldots \circ \fracConst_{i_m'+1} \circ \fracConst_{\alpha_{j_1}} \circ \fracConst_{\alpha_{j_2}} \circ \ldots)(K)\right) \\
		&= \mu \left((\fracConst_{i_1} \circ \ldots \circ \fracConst_{i_m+1})(K)\right) + \mu \left((\fracConst_{i_1} \circ \ldots \circ \fracConst_{i_m'+1})(K)\right)\\
		&= 2\cd 5^{-(m+1)}
\end{align}
and the result follows.
\end{proof}
\section{The Dirichlet Laplacian}\label{constructDirichletLaplacian}
We will now construct the Dirichlet Laplacian. Since the Dirichlet Laplacian is a densely defined operator we will therefore start by recalling some definitions for a densely defined operator on a Hilbert space. We will also need Dirichlet forms, thus we will also be defining them.

\medskip
\begin{defi}\label{selfAdjoint}
	Let $\mathcal{H}$ be a Hilbert space and let $A\colon \dom A \to \mathcal{H}$ be a densely defined operator. Then we say the following about $A$:
	\begin{enumerate}[(i)]
		\item $A$ is said to be symmetric if for $f,g \in \dom A$
		\begin{equation}
			\inn{Af}{g} = \inn{f}{Ag}.
		\end{equation}
		\item $A$ is said to be non-negative (non-positive resp.) if for $f \in \dom A$
		\begin{equation}
			\inn{Af}{f} \geq 0, \quad (\inn{Af}{f} \leq 0 \text{resp.}).
		\end{equation}
		\item $A$ is said to be self-adjoint if $\dom A^{*} = \dom A$ where $A^{*}$ is the adjoint of $A$ on the domain
		\begin{equation}
			\dom A^{*} = \cbr{f \in \mathcal{H} : \exists c(f) \geq 0, \forall g \in \mathcal{D}(A), \abs{\inn{f}{Ag}} \leq c(f) \norm{g}}
		\end{equation}
		where $c(f)$ denotes a constant that may depend on $f$. We can further define $A^*$ by the formula
		\begin{equation}
			\inn{A^*f}{g} = \inn{f}{Ag} , \quad \text{for } f \in \dom A^*, \  g \in \dom A.
		\end{equation}
	\end{enumerate}
\end{defi}

 \medskip
\begin{defi} A \emph{quadratic form} $\mathcal{E}$ on $\mathcal{H}$ is a non-negative, definite, symmetric bilinear form from $\dom \mathcal{E} \times \dom \mathcal{E} \to \ree$, where $\dom \mathcal{E}$ is a dense subspace of $\mathcal{H}$, where
	
	\begin{enumerate}[(i)]		
		\item A quadratic form is said to be closed if $\dom \mathcal{E}$ equipped with the norm
		\begin{equation}
			\norm{f}^{2}_{\dom \mathcal{E}} = \norm{f}^2 + \mathcal{E}(f,f)
		\end{equation}
		is a Hilbert space.
		\item A quadratic form $\mathcal{E}$ on $\mathcal{H}$ is said to be closable if it admits a closed extension. That is, there exists a closed quadratic form $\overline{\mathcal{E}}$ such $\dom \mathcal{E} \subset \dom \overline{\mathcal{E}}$ and $\overline{\mathcal{E}}$ coincides with $\mathcal{E}$ on $\dom \mathcal{E} \times \dom \mathcal{E}$.
		\item  We further call a quadratic form $\mathcal{E}$ a Dirichlet form if it has the following property. If $u, v \in \dom \mathcal{E}$ so that for almost every $x,y \in K$ we have
		\begin{equation}\label{contraction}
		\abs{v(x) - v(y)} \leq \abs{u(x) - u(y)}, \quad \text{and} \quad \abs{v(x)}, \leq \abs{u(x)}
		\end{equation}
		then 
		\begin{equation}
			\mathcal{E}(v,v) \leq \mathcal{E}(u,u).
		\end{equation}
	\end{enumerate} 
\end{defi} 
 
 
 
\noi Now recall from \Cref{renormalizedGraphEnergiesLimit} that 
\begin{equation}
	\mathcal{E}(f,f) = \lim_{m \to \infty} \mathcal{E}_m(f,f)
\end{equation}
with domain $\dom_{0} \mathcal{E}$ where
\begin{equation}
	\mathcal{E}_{m}(f,f) =  \sum_{x \underset{m}{\sim} y} (f(x)-f(y))^2.
\end{equation}

By \cite[Theorem 4.1]{specAnal}  we have that $(\mathcal{E}, \dom_{0}\mathcal{E})$ is a Dirichlet form on $L^{2}(K, \mu)$.  By \cite[Lemma 1.3.3]{fabriceDirichelt} we have that 
\begin{equation}
	\mathcal{E}(f,f) = - \int_{K}\Delta f g \dd \mu,
\end{equation}
where $\Delta$ a unique densely defined self-adjoint operator. The operator $\Delta$ is called the generator of $\mathcal{E}$. Moreover we have that $\Delta$ is an extension of the Kigami Laplacian $\Delta_{\mu}$. The operator $\Delta$ with domain $\dom \Delta$ is referred to as the Dirichlet Laplaican. 

Notice that the Kigami Laplacian $\Delta_{\mu}$ was not used in the construction of the Dirichlet Laplacian $\Delta$, but that we used the limit of the renormalized graph energy from \Cref{renormalizedGraphEnergiesLimit}. This stems from the fact that there are two ways to define $\Delta$, one through Dirichlet forms and one trough $\Delta_{\mu}$. We have chosen to define $\Delta$ through Dirichlet forms, but defining $\Delta$ through $\Delta_{\mu}$ follows a similarly condensed argument. One would prove that $\Delta_{\mu}\colon \dom \Delta_\mu \to L^{2}(K, \mu)$ is a densely defined non-positive symmetric operator. It then follows from \cite[Lemma 1.3.5]{fabriceDirichelt} that the quadratic form $\mathcal{E}(f,g) = -\inn{f}{\Delta_{\mu} g}$ for $f,g \in \dom \Delta_{\mu}$ is closable and the generator of $\mathcal{E}$ is the self-adjoint extension $\Delta$.

\medskip
\section{Discrete Fractional Laplacian}
Recalling \Cref{discreteMatrix} we consider the eigenvalues $0 < \lambda_1^{m} \leq \lambda_{2}^{m} \leq \ldots \leq \lambda_{N_m}^{m}$, each being repeated accordingly to multiplicity, of $-\Delta_m$. Next let $(\phi_i^{m})_{1 \leq i \leq N_m}$ be the corresponding eigenfunctions with respect to the measure $\mu_m$. This leads us to defining the following.

\medskip
\begin{defi}
	Let $s \geq 0$, we then define the \emph{discrete Riesz kernel} on $V_m$ as
	\begin{equation}
		G_s^{m}(x,y) = \sum_{i = 1}^{N_m} (\lambda_i^{m})^{-s} \phi_{i}^{m}(x)\phi_{i}^{m}(y), \quad \text{for }x, y \in V_m.
	\end{equation}
\end{defi}

By construction we can see that the matrix $\left(G_s^{m}(x,y)\right)_{x, y \in V_m}$ is symmetric. It is also positive semidefinite, which follows from
\begin{align}
	\sum_{i,j = 1}^{N_m}a_i a_j G_s^{m}(x_i, x_j) &= \sum_{i,j = 1}^{N_m}a_i a_j \sum_{k = 1}^{N_m} (\lambda_k^{m})^{-s} \phi_{k}^{m}(x_i)\phi_{k}^{m}(x_j) \\
	&=\sum_{k =1}^{N_m} (\lambda_k^{m})^{-s}  \left(\sum_{i = 1}^{N_m} \phi_k (x_i)^2 a_i^2 + \sum_{i, j = 1, i \neq j}^{N_m}\phi_k (x_i)\phi_k(x_j) a_ia_j\right) \\
	&= \sum_{k =1}^{N_m} (\lambda_k^{m})^{-s}  \left(\sum_{i = 1}^{N_m} \phi_k (x_i) a_i \right)^2 \geq 0.
\end{align}
Thus it can be the covariance matrix of a Gaussian vector. The Riesz kernel will allow us to define the discrete fractional Laplacian.
\begin{defi}\label{discreteFractionalLaplacian}
	For $f \in \ell(V_m)$ and $s \geq 0$, the \emph{discrete fractional Laplacian} \nl $(-\Delta_m)^{s}\colon L^{2}(V_m, \mu_m) \to L^{2}(V_m, \mu_m)$ is defined by
	\begin{equation}
		(-\Delta_{m})^{-s}f(x) = \frac{1}{a_m} \sum_{y \in V_m} G_{s}^{m}(x,y)f(y) = \sum_{i = 1}^{N_m} (\lambda_{i}^{m})^{-s} \phi_{i}^{m}(x) \frac{1}{a_m} \sum_{y \in V_m} \phi_{i}^{m}(y)f(y).
	\end{equation}
\end{defi}
\begin{remark}
	Note that $\frac{1}{a_m} \sum_{y \in V_m} \phi_{i}^{m}(y)f(y) = \langle\phi^{m}_i,f \rangle_{L^{2}(V_m, \mu_m)}$ and therefore that \nl $(-\Delta_{m})^{-s}f(x) = \sum_{i = 1}^{N_m} (\lambda_{i}^{m})^{-s} \phi_{i}^{m}(x) \langle \phi_i^{m},f \rangle_{L^2(V_m, \mu_m)}$.
\end{remark}

Having stated the discrete fractional Laplacian we move on to proving some of its useful properties.

\medskip
\begin{prop}\label{discreteFractionalLaplacianProperties}Let $(-\Delta_{m})^{-s}$ be defined as in \Cref{discreteFractionalLaplacian} then the following holds
	\begin{enumerate}[(i)]
			\item $(-\Delta_{m})^{-s}$ is a uniformly bounded operator. \label{DFLuniformConvergent}
			\item $(-\Delta_{m})^{-s}$ is a symmetric operator. \label{DFLsymmetric}
	\end{enumerate}

\end{prop}
\begin{proof}
	Here $\norm{\cd} = \norm{\cd}_{L^2(V_m, \mu_m)}$ and $\inn{\cd}{\cd} = \inn{\cd}{\cd}_{L^2(V_m, \mu_m)}$. We start with (\ref{DFLuniformConvergent}). We have
	\begin{align}
		\norm{(-\Delta_m)^{-s}f(x)}^2 &\leq \norm{\sum_{i = 1}^{N_m} (\lambda_{i}^{m})^{-s} \phi_{i}^{m}(x) \frac{1}{a_m} \sum_{y \in V_m} \phi_{i}^{m}(y)f(y).} \\
		& \leq \sum_{i = 1}^{N_m} (\lambda_{i}^{m})^{-2s} \norm{\phi_{i}^{m}(x) \frac{1}{a_m} \sum_{y \in V_m} \phi_{i}^{m}(y)f(y)}^2 \\
		& \leq (\lambda_1^{m})^{-2s} \sum_{i =1}^{N_m} \abs{\inn{\phi_i}{f}}^2 \\
		& \leq (\lambda_1^{m})^{-2s} \norm{f}^2,
	\end{align}
	where the second to last inequality comes from the fact that $0 < \lambda_1 < \lambda_{2} < \ldots$. The uniform bound comes from the fact that $\inf_{m} \lambda_{1}^{m} > 0$. This follows from \cite[Lemma 5.2]{specAnal}, since if $\inf_{m} \lambda_{1}^{m} \leq 0$ then the series in the Lemma would not converge.
	
	Moving on to (\ref{DFLsymmetric}), we have that
	\begin{align}
		\inn{f}{(-\Delta_{m})^{-s}g} & = \int_{V_m} f (-\Delta_{m})^{-s}g \dd \mu_m(x) \\
		&= \int_{V_m} f(x) \left(\sum_{j = 1}^{N_m} (\lambda_{j}^{m})^{-s} \phi_{j}^{m}(x) \frac{1}{a_m} \sum_{y \in V_m} \phi_{j}^{m}(y)g(y)\right) \dd \mu_m(x) \\
		&= \sum_{j = 1}^{N_m} (\lambda_{j}^{m})^{-s} \frac{1}{a_m} \left(\sum_{y \in V_m} \phi_{j}^{m}(y)g(y) \right)\int_{V_m} \left(\sum_{i =1}^{N_m}\phi^{m}_i(x) \inn{\phi^{m}_i}{f} \right) \phi^{m}_j(x) \dd \mu_m(x) \\
		&=\sum_{j = 1}^{N_m} (\lambda_{j}^{m})^{-s} \frac{1}{a_m} \left(\sum_{y \in V_m} \phi_{j}^{m}(y)g(y)\right) \int_{V_m} \phi^{m}_j(x) \inn{\phi^{m}_j}{f} \phi^{m}_j(x) \dd \mu_m(x) \\
		& = \sum_{j = 1}^{N_m} \bigg[ \int_{V_m} \left((\lambda_j^{m})^{-s} \phi_j(x) \inn{\phi_j^{m}}{f}\right) \phi_j^{m}(x) \inn{\phi_j^{m}}{g} \dd \mu_m(x) \\
		&\phantom{a} \qquad \qquad + \int_{V_m} \left((\lambda^{m}_{j})^{-s} \phi_j^{m}(x)\inn{\phi_j^{m}}{f} \right) \left(\sum_{i = 1, i \neq j}^{N_m} \phi_i^{m}(x) \inn{\phi_i^{m}}{g}\right) \dd \mu_m(x)\bigg] \\
		&= \int_{V_m} (-\Delta_{m})^{-s}f g \dd \mu_m(x) \\
		&=\inn{(-\Delta_{m})^{-s}f}{g}
	\end{align}
	using that $f = \sum_{i =1}^{N_m}\phi_i \inn{\phi_i}{f}$ and $\int_{V_m}\phi_i \phi_j \dd \mu_m$ is $1$ if $j = i$ and $0$ if $i \neq j$.
\end{proof}


\section{Semigroups and Heat kernels}
When constructing the fractional Laplacian in \Cref{sec:FractionalLaplacian}, we will use that $-\Delta$ generates a semigroup which admits a heat kernel. We will therefore first define a semigroup as well as showing that $-\Delta$ generates a semigroup. We will also define the heat kernel.

\begin{defi}\label{semigroupDef}
	Let $\mathcal{H}$ be a Hilbert space. A \emph{strongly continuous self-adjoint contraction semigroup} is a family of self-adjoint operators $(P_t)_{t \geq 0}\colon \mathcal{H} \to \mathcal{H}$ which satisfies the following properties:
	\begin{enumerate}[(i)]
		\item For $s,t \geq 0$, $P_t \circ P_s = P_{s +t}$ (semigroup property). \label{semigroup}
		\item For every $f \in \mathcal{H}, \lim_{t \to 0} P_t f= f$ (strong continuity).\label{strong cont}
		\item For every $f \in \mathcal{H}$ and $t \geq 0$, $\norm{P_t f} \leq \norm{f}$ (contraction property). \label{semigroup contraction}
	\end{enumerate}
\end{defi}

\medskip
\noi The following theorem is vital as it will show that an operator can generate a semigroup. 

\medskip
\begin{theorem}\label{genSemigroup}
	If $A$ is a densely defined self adjoint operator on $\mathcal{H}$ then it is the generator of the strongly continuous self-adjoint contraction semigroup on $\mathcal{H}$ defined as $P_t = e^{-tA}$.
\end{theorem}

\clearpage
\begin{proof}
	
	Let $A$ be given as stated. From the spectral theorem \cite[Theorem 1.1.2]{fabriceDirichelt} there is a measure space $(\Omega, \nu)$, a unitary map $U\colon L^2(\Omega, \nu) \to \mathcal{H}$ and a non negative real valued measurable function $\lambda$ on $\Omega$ such that
	\begin{equation}
		U^{-1}AUf(x) = -\lambda(x)f(x)
	\end{equation}
	for $x \in \Omega, U f \in \dom A$. Then from the functional calculus \cite[Theorem 1.1.3]{fabriceDirichelt} we define \nl $P_t\colon \mathcal{H} \to \mathcal{H}$ such that
	\begin{equation}
				U^{-1}P_tUf(x) = e^{-t \lambda(x)}f(x).
	\end{equation}
	Note that since $U$ is a unitary map, it is a linear bijection that for all $x,y \in L^2(\Omega, \nu)$ satisfies $\inn{Ux}{Uy}_{\mathcal{H}} = \inn{x}{y}_{L^2(\Omega, \nu)}$.
	
	It is enough to show that $Q_t = U^{-1}P_t U$ is self-adjoint, since then $P_t$ will be self-adjoint. This can be proven quite easily. If we have the functions $f, g \in \mathcal{H}$ and $\psi, \phi \in L^{2}(\Omega, \nu)$ then $P_t$ is symmetric since
	\begin{multline}
	\inn{Af}{g}_{\mathcal{H}} = \inn{AU\psi}{U\phi}_{\mathcal{H}} = \inn{UU^{-1}UAU\psi}{UU^{-1}U\phi}_{\mathcal{H}} = \inn{U^{-1}AU \psi}{\phi}_{L^2(\Omega, \nu)} \\
	=  \inn{\psi}{U^{-1}AU \phi}_{L^2(\Omega, \nu)} = \inn{U\psi}{UU^{-1}AU\phi}_{\mathcal{H}} = \inn{f}{Ag}_{\mathcal{H}}.
	\end{multline}
	That $\dom P_t^{*} = \dom P_t$ reasons from the same argument as in \eqref{operatorAndAdjointHaveSameDom}, which will be given later in the proof.
	
	We will prove that $Q_t = U^{-1}P_tU$ with domain $\dom Q_t$ is a self-adjoint strongly continuous contraction semigroup on $L^2(\Omega, \nu)$ with generator $A$. Starting with $A$ being a generator, we see that for $f \in \dom A$
	\begin{align}
		\lim_{t \to 0} \norm{\frac{Q_tf-f}{t} - Af} &= \lim_{t \to 0}\sqrt{\inn{\frac{Q_tf-f}{t} - Af}{\frac{Q_tf-f}{t} - Af}} \\
		&= \lim_{t \to 0}\sqrt{\inn{\frac{Q_tf-f}{t}}{\frac{Q_tf-f}{t} - Af} - \inn{Af}{\frac{Q_tf-f}{t} - Af}} \\
		&\leq \lim_{t \to 0}\sqrt{\inn{\frac{Q_tf-f}{t}}{\frac{Q_tf-f}{t} - Af}} \\
		&= \lim_{t \to 0}\sqrt{\int_{\Omega}\left(\frac{Q_tf-f}{t}\right)\left(\frac{Q_tf-f}{t} - Af\right)  \dd \nu} \\
		&= \lim_{t \to 0}\sqrt{\int_{\Omega}\left(\frac{(e^{-\lambda t}-1)f}{t}\right)\left(\frac{Q_tf-f}{t} - Af\right)} \\
			&= 0
	\end{align}
	where the convergence comes from the fact that $e^{-\lambda t}$ converges to $1$ much faster than $1/t$ diverges to $\infty$.  Thus $A$ is a generator.
	Next we have that $Q_t$ is self-adjoint, for which we start with symmetry
	\begin{equation}
		\inn{Q_t f}{g} = \int_{\Omega} Q_t f g \dd \nu = \int_{\Omega} e^{-\lambda t} f g \dd \nu = \inn{f}{Q_t g}.
	\end{equation}
	Furthermore $\dom Q_t$ is dense in $L^2(\Omega, \nu)$ since if $f_n \to f$ in $L^2(\Omega, \nu)$ and $Q_t f_n \to g$ then
	\begin{equation}
		\lim_{n \to \infty} Q_t f_n = \lim_{n \to \infty} e^{-\lambda t}f_n = e^{-\lambda t}f = Q_t f.
	\end{equation} 
	
	Showing that $\dom Q_{t}^{*} = \dom Q_t$. Since $Q_t$ is symmetric we have from the Cauchy-Schwartz inequality that
	\begin{equation}\label{operatorAndAdjointHaveSameDom}
		\abs{\inn{f}{Q_tg}} = \abs{\inn{Q_t f}{g}} \leq \norm{Q_t f} \norm{g}.
	\end{equation}
	Thus we have $Q_t$ is self-adjoint. 
	
	To show the semigroup properties, we have \Cref{semigroupDef}.(\ref{semigroup})
	\begin{equation}
		Q_{t+s}f = e^{-\lambda (t +s )}f = e^{-\lambda t}e^{-\lambda s}f = Q_t Q_s f.
	\end{equation}
	Further for \Cref{semigroupDef}.(\ref{strong cont}) we have
	\begin{equation}
		\lim_{t \to \infty} Q_t f = \lim_{t \to \infty} e^{-t\lambda} f = f
	\end{equation}
	and lastly \Cref{semigroupDef}.(\ref{semigroup contraction})
	\begin{equation}
		\norm{Q_t f} = \int_{\Omega} (e^{-\lambda t}f)^2 \dd \nu\leq \int_{\Omega} f^2 \dd \nu = \norm{f} .
	\end{equation}
	Thus $Q_t$ is a strongly continuous self-adjoint contraction semigroup on $L^2(\Omega, \nu)$ with generator $A$.
\end{proof}

The next corollary follows quite easily from \Cref{genSemigroup}.
\begin{coro}\label{coro:PtIsLinear}
	$P_t$ is a linear semigroup.
\end{coro}
\begin{proof}
	Let $f,g \in \mathcal{H}$ and $\alpha$ be some constant. Then 
	\begin{multline}
		U^{-1}P_tU\left(\alpha f(x) + g(x)\right) = e^{-t \lambda(x)}\left(\alpha f(x) + g(x)\right) \\
		= \alpha e^{-t \lambda(x)}f(x) + e^{-t \lambda(x)}g(x) = \alpha U^{-1}P_tU f(x) + U^{-1}P_tU g(x)
	\end{multline}
\end{proof}

With semigroups out of the way we can define the heat kernel.

\medskip
\begin{defi}\label{heatKernel}
	Let $(P_t)_{t \geq 0}$ be a strongly continuous self-adjoint contraction semigroup on $L^2(X, \mu)$. We say that the semigroup $(P_t)_{t \geq 0}$ \emph{admits a heat kernel} if the heat kernel measure has a density with respect to $\mu$, i.e. there exists a measurable function $p\colon \ree_{> 0} \times X \times X \to \ree_{\geq 0}$ such that for every $t > 0$, a.e. we have
	\begin{equation}
		P_tf(x) = \int_{K}p_t(x,y)f(y) \dd \mu(y)
	\end{equation}
	with $x,y \in X, f \in L^{\infty}(X, \mu)$.
\end{defi}

\medskip
\begin{remark}\label{remark:DiscreteLaplacianSemiGroup}
	The discrete Laplacian $-\Delta_m$ also generates a semigroup. This is because $\Delta_m$ can be represented as the matrix $L_m$ from \Cref{discreteMatrix}. Since $L_m$ is symmetric, it can be diagonalized in an orthonormal basis $(\Phi_i)_{1 \leq i \leq N_m}$, where $(\Phi_i)_j = \phi_i(x_j)$ for $x_i \in V_m$. We then have for every vector $X \in \ree^{N_m}$ that
	\begin{equation}
		e^{-tL_m}X = \sum_{i =1 }^{N_m} e^{-t \lambda_i}\inn{X}{\Phi_i}\Phi_i.
	\end{equation}
 	It is obvious to see that $P_t^{m} = e^{-tL_m}$ is a semigroup, as the reasoning follows the same given in \Cref{genSemigroup}. 	As we will want to use this semigroup on functions from $\ell(V_m)$ we will, for an $f \in \ell(V_m)$ and $x_j \in V_m$ write
	\begin{equation}
		e^{-tL_m}f(x_j) = \sum_{i=1}^{N_m} e^{-t \lambda_i} \inn{f}{\Phi_i}_{L^{2}(V_m, \mu_m)}\left(\Phi_i\right)_{j}
	\end{equation}
\end{remark}

\section{Fractional Laplacian}\label{sec:FractionalLaplacian}
Before we can define the continuous fractional Laplacian, we need to define the eigenvalues and eigenfunctions for $-\Delta$. 

Recall that the Laplacian $\Delta$, with domain $\dom\Delta$, is a densely defined self-adjoint operator. Therefore by \Cref{genSemigroup} it is the generator of a strongly continuous Markov semigroup $\cbr{P_t}_{t \geq 0}$ on $L^2(K, \mu)$. 

From \cite[Theorem 8.4]{barlow} this semigroup admits a heat kernel $p_t(x,y), t > 0, x,y \in K$, with respect to the normalized Hausdorff measure $\mu$. The heat kernel $p_t(x,y)$ is called the Dirichlet heat kernel on $K$. Furthermore we get that from \cite[section 2.3]{Fabrice} this heat kernel satisfies, for some $c_1, c_2 \in (0, \infty)$, 
\begin{equation}\label{heatKernelBound}
	p_t(x,y) \leq c_1 t^{-\sfrac{d_h}{d_w}} \exp\left(-c_2 \left(\frac{d(x,y)^{d_w}}{t}\right)^{\frac{1}{d_w -1}}\right)
\end{equation}
for every $x,y \in K \times K$ and $t \in (0, \infty)$ where $d_h$ is the Hausdorff dimension and $d_w$ is the walk dimension. 

\medskip
The Dirichlet heat kernel is quite central in the definition of the continuous fractional Laplacian, as its bounds and definition will be used quite a bit throughout the proofs in the following section. One of the more important properties of the heat kernel is, that it can be described as a sum of eigenvalues and eigenfunctions of $\Delta$. The following theorem gives us the existence of the eigenvalues, functions and the sum representation of the Dirichlet heat kernel. This theorem will not be proved, but a detailed proof can be found in \cite[Theorem 2.2]{Jozef}.

\medskip
\begin{theorem}\label{heatKernelWritten}
	There exists an orthonormal basis $\cbr{\phi_{i}}^{\infty}_{i = 1}$ of $L^2(K, \mu)$ with $\phi_j$ being eigenfunctions of $-\Delta$. If $\lambda_i$ is the eigenvalue belonging to $\phi_i$, then $\lambda_i >0$ and $\lim_{i \to \infty} \lambda_i = \infty$. Moreover
	\begin{equation}
		p_t(x,y) = \sum_{j = 1}^{\infty} e^{-\lambda_j t} \phi_j(x)\phi_j(y),
	\end{equation}
	where the convergence is uniform for all $t \in (0, \infty)$ and $x,y \in K \times K$.
\end{theorem}

The eigenvalues are not only used to represent the heat kernel but will also be used to represent the fractional Laplacian. The following lemma gives us an idea of what parameter value is need for convergence.
\begin{lemma}\label{convergenceOfEigenvaluesOfLaplacian}
	Let $0< \lambda_1 \leq \lambda_2 < \ldots$ be the eigenvalues of $-\Delta$, then the following holds
	\begin{equation}
		\sum_{j = 1}^{\infty} \frac{1}{\lambda_{j}^{2s}} < \infty \iff s > \frac{d_h}{2d_w}.
	\end{equation}
\end{lemma}
\begin{proof}
	This follows from \cite[Lemma 5.1.3]{Kigami_2001} which states that there exists constants $c_1, c_2 >0$ such that 
	\begin{equation} \label{convergenceOfEigenvaluesOfLaplcianBounds}
			c_1\cd n^{\sfrac{2}{d_s}} \leq \lambda_{n} \leq c_2\cd n^{\sfrac{2}{d_s}}.
	\end{equation}

	\noi Assume $s > \frac{d_h}{2d_w}$, then by \eqref{convergenceOfEigenvaluesOfLaplcianBounds}
	\begin{equation}
		\sum_{j = 1}^{\infty} \frac{1}{\lambda_j^{2s}} \leq \sum_{j = 1}^{\infty} \frac{1}{(c_1 n^{\sfrac{2}{d_s}})^{2s}} < \infty,
	\end{equation}
	since 
	\begin{equation}
		\frac{4s}{d_s} > \frac{\frac{2d_h}{d_w}}{\sfrac{2d_h}{d_w}} = 1.
	\end{equation}
	Hence the series is convergent.
	
	Assume now that 
	\begin{equation}
		\sum_{j = 1}^{\infty} \frac{1}{\lambda_{j}^{2s}} < \infty.
	\end{equation}
	Then again by \eqref{convergenceOfEigenvaluesOfLaplcianBounds} we have that
	\begin{equation}
						\sum_{j = 1}^{\infty} \frac{1}{(c_2 n^{\sfrac{2}{d_s}})^{2s}} \leq \sum_{j = 1}^{\infty} \frac{1}{\lambda_{j}^{2s}} < \infty,
	\end{equation}
	and the left sum only converges if $\sfrac{4s}{d_s} > 1$ in which case $s > \sfrac{d_s}{4} = \sfrac{d_h}{2d_w}$ as desired.
\end{proof}

Thus we are ready to define the fractional Laplacian. 

\medskip
\begin{defi}\label{fracLaplacian}
	Let $s \geq 0$, and $f \in L^2(K, \mu)$. The \emph{fractional Laplacian} $(-\Delta)^{-s}$ is defined as
	\begin{equation}
		(-\Delta)^{-s}f = \sum_{j = 1}^{\infty} \frac{1}{\lambda_{j}^{s}} \phi_j \int_{K} \phi_j(y)f(y) \dd \mu(y).
	\end{equation}
\end{defi}
\begin{remark}
	Once again we notice that $\inn{\phi_i}{f} = \int_{K} \phi_i f \dd \mu$, and can rewrite
	\begin{equation}
		(-\Delta)^{-s} f = \sum_{j = 1}^{\infty} \frac{1}{\lambda_j^{s}} \phi_j \inn{\phi_j}{f}.
	\end{equation}
\end{remark}

\clearpage
Now that we have the fractional Laplacian defined, we will show some of its very useful properties.
\medskip
\begin{prop}\label{fracLaplacianProperties}
	Let $(-\Delta)^{-s}$ be defined as in \Cref{fracLaplacian}. Then:
	\begin{enumerate}[(i)]
		\item 	$(-\Delta)^{-s}$ is a linear operator.\label{fracGausItem1}
		\item If $\phi_n$ is an orthonormal eigenfunction, then $(-\Delta)^{-s}\phi_n = \frac{1}{\lambda_{n}^{-s}}\phi_{n}$. \label{fracGausItem2}
		\item 	$(-\Delta)^{-s}$ is a bounded operator. \label{fracGausBounded}
		\item $(-\Delta)^{-s}$ is a symmetric operator. \label{fracGausSymmetric}
	\end{enumerate}
\end{prop}
\begin{proof}
	
	We begin with the proof of (\ref{fracGausItem1}).
	
	\noi Let $f,g \in L^{2}(K, \mu)$ and let $\alpha$ be some constant. Then by definition of $(-\Delta)^{-s}$ we have
	\begin{align}
		(-\Delta)^{-s}(\alpha f + g) &= \sum_{j = 1}^{\infty} \frac{1}{\lambda_{j}^{s}} \phi_j \int_{K} \phi_j(y)(\alpha f(y) + g(y)) \dd \mu(y)\\
		& = \alpha \sum_{j = 1}^{\infty} \frac{1}{\lambda_{j}^{s}} \phi_j \int_{K} \phi_j(y)f(y)  \dd \mu(y) + \sum_{j = 1}^{\infty} \frac{1}{\lambda_{j}^{s}} \phi_j \int_{K} \phi_j(y)g(y) \dd \mu(y) \\
		&=\alpha (-\Delta)^{-s}f + (-\Delta)^{-s}g.
	\end{align}
	 Moving on to (\ref{fracGausItem2}). We note that
	\begin{equation}
		\int_{K}\phi_j \phi_i \dd \mu = \delta_{ij},
	\end{equation}
	where $\delta_{ij}$ is the Kronecker delta. We then have for any eigenfunction $\phi_n$ that
	\begin{equation}
		(-\Delta)^{-s}\phi_n = \sum_{j = 1}^{\infty} \frac{1}{\lambda_{j}^{s}} \phi_j(x) \int_{K}\phi_j(y)\phi_n(y) \dd \mu(y) = \frac{1}{\lambda_{n}^{-s}}\phi_{n}.
	\end{equation}
	Next we have (\ref{fracGausBounded}). This follows from
	\begin{align}
		\norm{(-\Delta)^{-s}f}^2 &= \norm{\sum_{j = 1}^{\infty} \frac{1}{\lambda_{j}^{s}} \phi_j \inn{\phi_j}{f}}^2 \\
			&= \left\langle \sum_{j = 1}^{\infty} \frac{1}{\lambda_{j}^{s}} \phi_j \inn{\phi_j}{f},\sum_{j = 1}^{\infty} \frac{1}{\lambda_{j}^{s}} \phi_j \inn{\phi_j}{f} \right\rangle \\
			&= \sum_{j=1}^{\infty} \abs{\frac{1}{\lambda_j^s} \inn{\phi_j}{f}}^2\inn{\phi_j}{\phi_j} \\
			&\leq \left(\frac{1}{\lambda_1^s}\right)^2 \sum_{j=1}^{\infty} \abs{\inn{\phi_j}{f}}^2 \\
			&= \left(\frac{1}{\lambda_1^s}\right)^2 \norm{f}^2,
	\end{align}
	in which case we have our bound. Finally we will show (\ref{fracGausSymmetric})
	\begin{align}
		\inn{(-\Delta)^{-s}f}{g} &= \int_{K} \left(\sum_{j = 1}^{\infty} \frac{1}{\lambda_{j}^{s}} \phi_j \inn{\phi_j}{f} \right) g \dd \mu \\
		&=\int_{K} \left(\sum_{j = 1}^{\infty} \frac{1}{\lambda_{j}^{s}} \phi_j \inn{\phi_j}{f} \right) \left(\sum_{i = 1}^{\infty} \phi_i \inn{\phi_i}{g}\right) \dd \mu \\
		&= \sum_{j = 1}^{\infty} \int_{K} \frac{1}{\lambda_j^s}\phi_j \inn{\phi_j}{f}\left(\sum_{i = 1}^{\infty} \phi_i \inn{\phi_i}{g}\right) \dd \mu \\
		&= \sum_{j =1}^{\infty} \int_{K} \frac{1}{\lambda_j^{s}} \phi_j^2 \inn{\phi_j}{f}\inn{\phi_j}{g} \dd \mu \\
		&= \sum_{j = 1}^{\infty} \left[\int_{K} \frac{1}{\lambda_j^{s}} \phi_j \inn{\phi_j}{f} \phi_j\inn{\phi_j}{g} \dd \mu + \int_{K}\frac{1}{\lambda_j^{s}}\phi_j\inn{\phi_j}{g} \left(\sum_{i = 1, i \neq j}^{\infty} \phi_i \inn{\phi_i}{f}\right) \dd \mu\right] \\
		&= \sum_{j=1}^{\infty} \int_{K}\left(\frac{1}{\lambda_j^{s}}\phi_j\inn{\phi_j}{g}\right)\left(\sum_{i = 1}^{\infty} \phi_i \inn{\phi_i}{f}\right) \dd \mu \\
		& = \int_{K} \left(\sum_{j = 1}^{\infty} \frac{1}{\lambda_j^s} \phi_j \inn{\phi_j}{g}\right) f \dd \mu \\
		&= \inn{f}{(-\Delta)^sg)}
	\end{align}
	as desired.
\end{proof}
While the definition of the fractional Laplacian is not the most complicated, we are able to represent $(-\Delta)^{-s}$ in another way, namely with the Riesz kernel.

\medskip
\begin{defi}\label{fractionalRieszKernel}
	For parameter $s \geq 0$ we define the \emph{fractional Riesz kernel} $G_s$ by
	\begin{equation}
		G_{s}(x,y) = \frac{1}{\Gamma(s)}\int_{0}^{\infty} t^{s-1}p_t(x,y) \dd t, \quad x, y \in K, x \neq y.
	\end{equation}
\end{defi}
\begin{remark}
	It follows from \eqref{heatKernelBound} that $G_s(x,y)$ converges for every $s \geq 0$ with $x \neq y$.
\end{remark}

\medskip
The Riesz kernel allows us to use the estimates of the heat kernel $p_t(x,y)$ in conjunction with $(-\Delta)^{-s}$ as we shall shortly see. We will start by showing some of the bounds of the fractional Riesz kernel. 
\clearpage
\begin{prop}\label{boundsOfRieszKernel}
	Let $G_s(x,y)$ be as in \Cref{fractionalRieszKernel} then the following bounds hold.
	\begin{enumerate}[(i)]
		\item If $s \in [0, d_h/d_w)$, there exists a constant $C > 0$ such that for every $x,y \in K, x \neq y$
		\begin{equation}
			G_s(x,y) \leq \frac{C}{d(x,y)^{d_h-sd_w}}.
		\end{equation}
		
		\item If $s = d_h/d_w$, there exists a constant $C > 0$ such that for every $x,y \in K, x \neq y$
			\begin{equation}
				G_s(x,y) \leq C \abs{\ln \left(d(x,y) \right)}.
			\end{equation}
			
		\item If $s > d_h/d_w$, there exists a constant $C > 0$ such that for every $x,y \in K$
		\begin{equation}
			G_s(x,y) \leq C.
		\end{equation}
	\end{enumerate}
\end{prop}
Before proving \Cref{boundsOfRieszKernel} we will quickly state and prove a lemma that is needed.

\medskip
\begin{lemma}\label{heatKernelBound2}
	There exists a constant $C> 0$ such that for every $x,y \in K$ and $t \geq 1$,
	\begin{equation}
		\abs{p_t(x,y)} \leq Ce^{-\lambda_1 t}
	\end{equation}
	where $\lambda_1 > 0$ is the first non-zero eigenvalue of $-\Delta$.
\end{lemma}
\begin{proof}
	We have from \Cref{heatKernelWritten} that $p_t(x,y)$ has the expansion
	\begin{equation}\label{returnHeatKernel}
		p_t(x,y) = \sum_{j = 1}^{\infty} e^{-\lambda_j t}\phi_j(x)\phi_j(y).
	\end{equation}
	Since $\phi_j$ are eigenvalues, we have from definition of $p_t$, that for any $t > 0$ we have
	\begin{equation}
		\phi_j(x) = e^{\lambda_jt}\int_{K}p_t(x,y)\phi_j(y) \dd \mu (y).
	\end{equation}
	Then, from the Cauchy-Schwartz inequality, we have for every $t > 0$ 
	\begin{equation}
		\abs{\phi_j(x)} \leq e^{\lambda_j t} \left(\int_{K} p_t(x,y)^{2} \dd \mu (y)\right)^{\sfrac{1}{2}}\left(\int_{K} \phi_j(y)^{2} \dd \mu (y)\right)^{\sfrac{1}{2}}  = e^{\lambda_j t}p_{2t}(x,x)^{\sfrac{1}{2}}.
	\end{equation}
	The last equality comes from the fact that $\int_{K}\phi_i(y)\phi_j(y) \dd \mu = \delta_{ij}$ so $\int_{K} \phi_j(y)^{2} \dd \mu (y) = 1$. That combined with \eqref{returnHeatKernel} gives us $\int_{K} p_t(x,y)^{2} \dd \mu (y) = \int_{K} p_{2t}(x,y) \dd \mu (y)$ resulting in the equality.
	
	Now choosing $t = 1/4$ together with \eqref{heatKernelBound}, we have that there is a constant $C > 0$ so that for every $x \in K$ we have
	\begin{equation}
		\abs{\phi_j} \leq C e^{\sfrac{\lambda_j}{4}}.
	\end{equation}
	Returning to \eqref{returnHeatKernel} we have for every $x, y \in K$ and $t \geq 1$ that
	\begin{align}
		\abs{p_t(x,y)} &\leq \sum_{j = 1}^{\infty} e^{-\lambda_jt} \abs{\phi_j(x)} \abs{\phi_j(x)} \\
		& \leq C^2 \sum_{j = 1}^{\infty} e^{-\lambda_jt} e^{\sfrac{\lambda_j}{2}} \\
		& = C^2 \sum_{j = 1}^{\infty} e^{-\lambda_jt} e^{\sfrac{\lambda_j}{2}} e^{\lambda_j}e^{-\lambda_j} \\
		&\leq C^2 e^{-\lambda_1 t}e^{\lambda_1}\sum_{j = 1}^{\infty} e^{-\sfrac{\lambda_j}{2}},
	\end{align}
	where the last inequality comes from the fact that $0 < \lambda_1 < \lambda_2 < \ldots$ and $t \geq 1$ so \\ $e^{\lambda_j(1-t)} \leq e^{\lambda_1(1-t)}$. Thus we have the desired bound.
\end{proof}

\begin{proof}[Proof of \Cref{boundsOfRieszKernel}]
	We have that
	\begin{align}
		G_s(x,y) &= \frac{1}{\Gamma(s)}\int_{0}^{\infty} t^{s-1}p_t(x,y) \dd t \\
		&=\frac{1}{\Gamma(s)}\int_{0}^{1} t^{s-1}p_t(x,y) \dd t + \frac{1}{\Gamma(s)}\int_{1}^{\infty} t^{s-1}p_t(x,y) \dd t.
	\end{align}
	The integral $\frac{1}{\Gamma(s)}\int_{1}^{\infty} t^{s-1}p_t(x,y) \dd t$ can be uniformly bounded by a constant using \Cref{heatKernelBound2}. Thus we only need to bound the integral $\int_{0}^{1} t^{s-1}p_t(x,y) \dd t$. By \eqref{heatKernelBound} we have that
	\begin{equation}
		\int_{0}^{1} t^{s-1}p_t(x,y) \dd t \leq c_1 \int_{0}^{1} t^{s-1-\sfrac{d_h}{d_w}} \exp\left(-c_2 \left(\frac{d(x,y)^{d_w}}{t}\right)^{\frac{1}{d_w -1}}\right) \dd t.
	\end{equation}
	Since the bound of the integral relies on the value of $s$ we will split it up into the cases of \Cref{boundsOfRieszKernel}.
	
	\medskip
	\noi If $s > d_h/d_w$ then we simply have 
	\begin{equation}
		\int_{0}^{1} t^{s-1-\sfrac{d_h}{d_w}} \exp\left(-c_2 \left(\frac{d(x,y)^{d_w}}{t}\right)^{\frac{1}{d_w -1}}\right) \dd t \leq \int_{0}^{1}t^{s-1-\sfrac{d_h}{d_w}} \dd t.
	\end{equation}
	Next, if $s < d_h/d_w$, then using a change of variable $t = u d(x,y)^{d_w}$, we find that
	\begin{align}
		&\phantom{=}\int_{0}^{1} t^{s-1-\sfrac{d_h}{d_w}} \exp\left(-c_2 \left(\frac{d(x,y)^{d_w}}{t}\right)^{\frac{1}{d_w -1}}\right) \dd t \\
		&=d(x,y)^{sd_w - d_h} \int_{0}^{\sfrac{1}{d(x,y)^{d_w}}} u^{s-1-\sfrac{d_h}{d_w}} \exp\left(-c_2 \left(\frac{1}{u}\right)^{\frac{1}{d_w -1}}\right) \dd u \\
		&\leq d(x,y)^{sd_w-d_h} \int_{0}^{\infty} u^{s-1-\sfrac{d_h}{d_w}} \exp\left(-c_2 \left(\frac{1}{u}\right)^{\frac{1}{d_w -1}}\right) \dd u.
	\end{align}
	Convergence of the last integral comes from the fact that $e^{-c_2(1/u)^{\sfrac{1}{d_w}}}$ converges faster than $1/u$ diverges.
	
	\noi Lastly for $s = d_h/d_w$, we again use the change of variable $t = ud(x,y)^{d_w}$, setting $R = \diam(K)$ and noting that $0 < 1/R^{d_w} < 1/d(x,y)^{d_w}$, we thus find
	\begin{align}
		&\phantom{=}\int_{0}^{1} t^{s-1-\sfrac{d_h}{d_w}} \exp\left(-c_2 \left(\frac{d(x,y)^{d_w}}{t}\right)^{\frac{1}{d_w -1}}\right) \dd t \\
		&=\int_{0}^{\sfrac{1}{d(x,y)^{d_w}}} u^{-1} \exp\left(-c_2 \left(\frac{1}{u}\right)^{\frac{1}{d_w -1}}\right) \dd u \\
		&\leq \int_{\sfrac{1}{R^{d_w}}}^{\sfrac{1}{d(x,y)^{d_w}}} u^{-1} \exp\left(-c_2 \left(\frac{1}{u}\right)^{\frac{1}{d_w -1}}\right) \dd u + \int_{0}^{\sfrac{1}{R^{d_w}}} u^{-1} \exp\left(-c_2 \left(\frac{1}{u}\right)^{\frac{1}{d_w -1}}\right) \dd u  \\
		&\leq \int_{\sfrac{1}{R^{d_w}}}^{\frac{1}{d(x,y)^{d_w}}} u^{-1} \dd u + \int_{0}^{\sfrac{1}{R^{d_w}}} u^{-1} \exp\left(-c_2\left(\frac{1}{u}\right)^{\frac{1}{d_w -1}}\right) \dd u  \\
		& \leq d_w \abs{\ln \left(d(x,y)\right)} +\int_{0}^{\sfrac{1}{R^{d_w}}} u^{-1} \exp\left(-c_2\left(\frac{1}{u}\right)^{\frac{1}{d_w -1}}\right) \dd u \\
		& \leq C \abs{\ln \left(d(x,y) \right)}.
	\end{align}
	Thus we have our desired bounds.
\end{proof}

\section{Bounds of the Fractional Laplacian}
In the following we will prove some bounds for the fractional Laplacian. Since the goal of this thesis is to prove convergence of fractional Gaussian fields within a specific space of test functions, we will only be proving the lemmas for that space. The space of functions is given by

\medskip
\begin{defi}
	We define the following space of test functions $S(K)$ as
	\begin{equation}
		S(K) = \cbr{f \in C_0(K) \colon \forall k \geq 0 \lim_{n \to \infty} n ^{k} \abs{\int_{K} \phi_n(y)f(y) \dd \mu (y)} = 0}
	\end{equation}
	with the topology defined by the family of norms
	\begin{equation}
		\norm{f}_k = \norm{(-\Delta)^{-k}f}_{L^2(K, \mu)}.
	\end{equation}
\end{defi}
By \cite[Theorem 2]{germanOne} and \Cref{convergenceOfEigenvaluesOfLaplacian} we have that $S(K)$ is a Frèchet nuclear space. That a space is \emph{Fréchet} means that it is nuclear and the topology is given by a countable family of seminorms. We will quickly define a nuclear space, but will not delve further into the subject matter. The following definition is from \cite[Chapter 50]{Treves}.

\medskip
\begin{defi} \label{NuclearSpace}
	Let $E$ be a locally convex Hausdorff space and $\norm{\cdot}_p$ a continuous seminorm on $E$. Furthermore  let $\widehat{E}_{p}$ be a Banach space, defined to be the completion of the normed space $E \setminus \ker \norm{\cdot}_p$. Then $E$ is a \emph{nuclear space} if for every continuous seminorm $\norm{\cdot}_p$ on $E$ there exists  another continuous seminorm on $E$ with $\norm{\cdot}_q \geq \norm{\cdot}_p$ such that the canonical mapping $\widehat{E}_{p} \to \widehat{E}_q$ is \emph{nuclear}. 
	
	A mapping $u\colon E \to F$, where $E,F$ are Banach spaces, is \emph{nuclear} if there is a sequence $\cbr{x'_k}$ in the closed unit ball of $E'$, a sequence $\cbr{y_k}$ in the closed unit ball of $F$, and a complex sequence $\cbr{\lambda_k}$ with $\sum_{k}\abs{\lambda_k} < \infty$ such that $u(x) = \sum_{i = 1}^{\infty} \lambda_i x'_i(x)y_i$ for all $x \in X$. 
	
	Furthermore, the trace-norm $\norm{u}_{\text{Tr}}$ is equal to the infimum of the numbers $\sum_{k} \abs{\lambda_k}$ over the set of all representations of $u$ as such a series.
\end{defi}
\begin{remark}
	We have just defined a nuclear map on a Banach spaces, but note that the definition can be generalized to locally convex Hausdorff spaces.
\end{remark}

\medskip
The concept of a nuclear space is quite complicated, but its properties are very useful. It will allow us to prove the convergence of fractional Gaussian fields through the characteristics function and existence of fractional Gaussian fields through Bochner-Minlos theorem for nuclear spaces, which will be state later in \Cref{BochnerMinlos}. But first we will prove some lemmas regarding the fractional Laplacian and the fractional Riesz kernel.

\medskip
\begin{lemma}\label{fractionalLaplacianIsRieszKernel}
	Let $s > 0$. For $f \in S(K)$ and $x \in K$
	\begin{equation}
		(-\Delta)^{-s}f = \int_{K}G_s(x,y)f(y) \dd \mu(y).
	\end{equation}
\end{lemma}
\begin{proof}
	First, we will make sure the integral on the right side actually converges. This holds because $S(K) \subset C(K)$ hence $f$ will be bounded on $K$. Thus we just need for every $x \in K$, that $\int_{K}G_s(x,y) \dd \mu(y)$ is convergent.\footnote{In \Cref{boundsOfRieszKernel}, (i) and (ii) the bounds are for $x \neq y$, but they are null-sets with respect to $\mu$}  For $s \in (d_h/d_w, \infty)$ the convergence follows from \cref{boundsOfRieszKernel}. Looking at $s \in (0, d_h/d_w]$, we will use a dyadic annuli decomposition to split up $K$ into infinitely many annuli on the form $ B(x, R2^{-j})\setminus B(x,RS^{-j-1})$, noting that
	\begin{equation}
		K \subseteq \bigcup_{j \geq 0} B(x, R2^{-j})\setminus B(x,RS^{-j-1}),
	\end{equation}
	where $R = \diam(K)$. For all $y \in B(x, R2^{-j})\setminus B(x,RS^{-j-1})$ we then have that
	\begin{equation}
		R2^{-j-1} \leq d(x,y) \leq R2^{-j-1} \Rightarrow 2^{-j-1} \leq R^{-1} \cd d(x,y) \leq 2^{-j-1} .
	\end{equation}
	  We will show that the integral is finite for both $s = d_h/d_w$ and $s<d_h/d_w$, but this requires that we bound $\abs{\ln \left(d(x,y) \right)}$. So we see that for all $\alpha > 0$
	\begin{equation}
		x^{\alpha} \abs{\ln \left(x\right)} \to 0, \quad \text{as } x \to 0,
	\end{equation}
	which means there exists a constant $k > 0$ such we can choose some interval $(0, u_k]$ so 
	$\ln(x) \leq \frac{k}{x^{\alpha}}$ for all $x \in (0, u_k]$. \clearpage \noi Now choose $k'$ such $u_{k'} \leq R$, in which case for $s = d_h/d_w$ we have
	\begin{align}
		\int_{K} G_s(x,y) \dd \mu(y) &\leq \int_{K} C_1 \cd \abs{\ln\left( d(x,y)\right)} \dd \mu(y)\\
		& \leq  \sum_{j = 0}^{\infty} \int_{B(x, R2^{-j})\setminus B(x,RS^{-j-1})} C_1\cd \abs{\ln \left(d(x,y)\right)} \dd \mu(y) \\
		&\leq \sum_{j = 0}^{\infty} \int_{B(x, R2^{-j})\setminus B(x,RS^{-j-1})} \frac{C_1 \cd k'}{d(x,y)^{\gamma}} \dd \mu,
	\end{align}
	where $\gamma < d_h$. We can likewise bound $G_{s}(x,y)$ by
	\begin{align}
		\int_{K} G_s(x,y) \dd \mu(y) &\leq \int_{K} \frac{C_2}{d(x,y)^{\gamma}}\dd \mu(y)\\
		& \leq  \sum_{j = 0}^{\infty} \int_{B(x, R2^{-j})\setminus B(x,RS^{-j-1})} \frac{C}{d(x,y)^{\gamma}}\dd \mu(y), 
	\end{align}
	where $C = \max\cbr{C_1 \cd k', C_2}$. Thus for $s \in (0, d_h/d_w]$ we have
	\begin{align}
		\int_{K} G_s(x,y) \dd \mu(y) & \leq  \sum_{j = 0}^{\infty} \int_{B(x, R2^{-j})\setminus B(x,RS^{-j-1})} \frac{C}{d(x,y)^{\gamma}}\dd \mu(y) \\
		& \leq R\cd C \sum_{j = 0}^{\infty} 2^{j \gamma} \mu \left(B(x, R2^{-j})\setminus B(x,RS^{-j-1})\right)  \\
		&\leq R\cd C \sum_{j = 0}^{\infty} 2^{j \gamma} \mu \left(B(x, R2^{-j})\right) \\
		& \leq R^{d_h +1} \cd C \cd c_2 \sum_{j = 0}^{\infty} 2^{j(\gamma - d_h)} < \infty,
	\end{align}
	where we use \cref{normalziedHausdorffMeasureAhlforsRegular}. Therefore $\int_{K}G_s(x,y)\dd \mu(y)$ is convergent for every $x \in K$.
	
	\medskip
	Moving along in the proof, we note that if we use the substitution $t = \lambda_j u$ we get 
	\begin{align}
		\frac{1}{\Gamma(s)} \int_{0}^{\infty} t^{s-1}e^{-\lambda_j t} \dd t = \frac{1}{\Gamma(s) \cd \lambda_j^s} \int_{0}^{\infty} u^{s-1}e^{-u} \dd u = \frac{1}{\Gamma(s)\lambda_{j}^{s}} \Gamma(s) = \lambda_j^{-s}.
	\end{align}
	
	\noi Now using Fubini's theorem, \Cref{fractionalRieszKernel} and \Cref{heatKernelWritten}  we get
	\begin{align}
		\int_{K} G_s(x,y) f(y) \dd \mu (y) & = \int_{K} \frac{1}{\Gamma(s)}\int_{0}^{\infty} t^{s-1}p_t(x,y) \dd t f(y)\dd \mu (y) \\
		&= \int_{K} \frac{1}{\Gamma(s)}\int_{0}^{\infty} t^{s-1} \sum_{j = 1}^{\infty} e^{-\lambda_j t} \phi_j(x)\phi_j(y)  \dd t f(y)\dd \mu (y) \\
		&= \sum_{j = 1}^{\infty} \phi_j(x) \int_{K}\phi_{j}(y) f(y)  \frac{1}{\Gamma(s)} \int_{0}^{\infty} t^{s-1}e^{-\lambda_j t} \dd t \dd \mu(y) \\
		&= \sum_{j = 1}^{\infty} \frac{1}{\lambda_{j}^{s}} \phi_j(x) \int_{K}\phi_j(y)f(y) \dd \mu(y) \\
		&= (-\Delta)^{-s}f
 	\end{align}
 	as desired.
\end{proof}

An immediate corollary follows:
\begin{coro}\label{coro:ForRieszKernelAndFractionalLaplacian} Let $(-\Delta)^{-s}$ be defined as in \Cref{fracLaplacian}, then the following holds
	\begin{enumerate}[(i)]
		\item 	Let $s > 0$, for $f,g \in S(K)$ \label{coroRiesz1}
		\begin{equation}
			\int_{K} (-\Delta)^{-s} f (-\Delta)^{-s} g \dd \mu = \int_{K}\int_{K} f(x)g(y)G_{2s}(x,y) \dd \mu(x) \dd \mu (y).
		\end{equation}
		
		\item Let $s > 0$ the for $f \in S(K)$\label{coroRiesz2}
		\begin{equation}
			(-\Delta)^{-s}f = \frac{1}{\Gamma(s)}\int_{0}^{\infty} t^{s-1}P_tf(x) \dd t.
		\end{equation}
	\end{enumerate}

\end{coro}
\begin{proof}
	Beginning with (\ref{coroRiesz1}). It follows from \Cref{fracLaplacianProperties} and \Cref{fractionalLaplacianIsRieszKernel} that
	\begin{align}
		\int_{K} (-\Delta)^{-s}f (-\Delta)^{-s}g \dd \mu &= \inn{(-\Delta)^{-s}f}{(-\Delta)^{-s}g} \\
		& = \inn{f}{(-\Delta)^{-2s}g}\\
		 & = \int_{K}f (-\Delta)^{-2s}g \dd \mu \\
		 & =\int_{K}\int_{K} f(x)g(y)G_{2s}(x,y) \dd \mu(x) \dd \mu (y).
	\end{align}
	
	Moving on to (\ref{coroRiesz2}). Again from \Cref{fractionalLaplacianIsRieszKernel} and \Cref{heatKernelWritten} we have
	\begin{align}
		(-\Delta)^{-s}f &= \sum_{j =1 }^{\infty} (\lambda_j)^{-s} \phi_j \int_{K}\phi_j(y) f(y) \dd \mu(y) \\
		&= \sum_{j =1 }^{\infty} \left(\frac{1}{\Gamma(s)} \int_{0}^{\infty} t^{s-1}e^{-\lambda_j t} \dd t\right) \phi_j \int_{K}\phi_j(y) f(y) \dd \mu(y) \\
		&= \frac{1}{\Gamma(s)} \int_{0}^{\infty} t^{s-1} \int_{K} \sum_{j = 1}^{\infty}  e^{-\lambda_j t}  \phi_j \phi_j(y) f(y) \dd \mu(y) \dd t \\
		&= \frac{1}{\Gamma(s)} \int_{0}^{\infty} t^{s-1} \int_{K} p_t(x,y) f(y) \dd \mu(y) \dd t \\
		&=\frac{1}{\Gamma(s)} \int_{0}^{\infty} t^{s-1} P_tf(x) \dd t.
	\end{align}
\end{proof}